\documentclass[11pt]{article}
\usepackage[a4paper,margin=30mm]{geometry}
\usepackage{amsmath,amssymb,amsthm,mathtools}
\usepackage{enumitem}
\usepackage[hidelinks]{hyperref}
\setlength{\emergencystretch}{2em}

\newtheorem{theorem}{Theorem}[section]
\newtheorem{proposition}[theorem]{Proposition}
\newtheorem{lemma}[theorem]{Lemma}
\newtheorem{corollary}[theorem]{Corollary}
\newtheorem{definition}[theorem]{Definition}
\newtheorem{remark}[theorem]{Remark}
\newcommand{\cL}{\mathcal L}
\newcommand{\cU}{\mathcal U}
\newcommand{\cD}{\mathcal D}
\newcommand{\cK}{\mathcal K}
\newcommand{\cB}{\mathcal B}
\newcommand{\cQ}{\mathcal Q}
\newcommand{\one}{\mathbf 1}
\newcommand{\TV}{\mathrm{TV}}
\newcommand{\Lip}{\mathrm{Lip}}

\title{The Two-Time Moving-Defect Problem for Dispersing Billiards\\
\large Branch-section finite rank, positive-envelope oscillation, and a conditional
pathwise CM2 criterion}
\author{Working note}
\date{13 July 2026 (version 30)}

\begin{document}
\maketitle

\begin{abstract}
We isolate the two-time estimate needed to differentiate an infinite
correlation susceptibility under a genuine motion of dispersing scatterers.
The moving first defect contains regular two-dimensional terms and singular
one-dimensional physical faces.  Neither a projective-cone trace estimate nor
a nuclear graph-current representation is available.  We avoid both.  In the
central two-time band we first remove contracted coordinates.  The physical
face is then pushed to its two scalar dominant coordinates, smoothed only on
that quotient, and lifted on two same-branch homogeneous sections.  The
branchwise contracted-coordinate estimate intertwines the original graph
pairing with the section pairing without averaging over other symbolic pasts
or singularity cuts.  Thus a two-dimensional Fourier polynomial is never
declared to be a four-dimensional density by a fictitious chart inverse.  At
finite cutoff the lifted low block is exactly a sum of products of measured
standard-curve currents.  Oscillatory section densities are decomposed into
positive regular standard families; independent forward and reverse
one-time coupling estimates give product decay with an explicit polynomial
dominant-frequency loss.  Current centering and its test-side adjoint remain
distinct typed operations.  The complementary dominant
modes are treated by the physical phase estimate; no ambient Fourier vector is
incorrectly identified with a tangent frequency.  Off the
central band a signed proper-family theorem supplies forward and reverse
coupling.  This removes graph Banach spaces, tensor renewal, and nuclear
decompositions of the regular defect.

The high-frequency face term is controlled by a physical slope energy.  We
use Leclerc's native unstable coordinate and his cutting theorem modulo every
fixed-degree polynomial.  A finite-depth thickening condition records the
robust child margin, finite-time holonomies, actual-SRB cylinder weights, and
an exact stopping-antichain cover.  Every good symbolic child is then refined
to a second physical scale on which the total Lipschitz variation of the
original quantity $X^{[J]}-P$ is smaller than the core margin.  No affine jet
is absorbed into the polynomial defining the event.  The construction never compares
an infinite fractal conditional with an absolutely continuous one.  A single-source
conditional identity freezes one oriented endpoint and places the other
endpoint in the context polynomial.  A bounded-overlap native stopping cover, a mass-preserving
resolving buffer, and Jensen give the reduced-slope small-ball estimate.  The
choice $\sigma=r^{1/(1+\alpha_L)}$ separates Leclerc's cutting exponent from
the billiard $C^{1+\alpha_B}$ remainder.  On the nonresonant set a direct
$C^{1+\alpha}$ fractional nonstationary-phase estimate yields a polynomial
Fourier saving, which becomes exponential after a cutoff
$K=e^{\eta\max\{m,n\}}$ in the balanced band.

For radial faces we index every connected component and normalize its mass
before forming the boundary functional.  A full tangent one-dimensional area
formula replaces the vertical two-flux formula, which is retained only as a
special case.  It is combined with an oriented component ledger; no component is silently
discarded.  The resulting standard families recover after $O(\log Z)$
iterates.  Qualitative residual transversality is insufficient for the
infinite-depth limit, so the theorem is stated for analytic or circular table
deformations satisfying an explicit pathwise jet-tail condition with a
strict return/grazing margin.  A depth $L(s)\asymp\log(1/|s|)$ and a second
time threshold $T(s)\asymp\log(1/|s|)$ are used: shallow faces remain
transverse, short-time deep pieces are controlled in total variation, and
long-time central collars are controlled by a boundary-crossing estimate for
one positive envelope that simultaneously dominates the actual multi-time law
and every functional-correlation split law, followed by one common
$N$-dependent cylinder approximation;
off-diagonal collars use moving standard-family coupling.  This gives an
$o(1)$ tail in the absolute CM2 double sum, not merely pointwise convergence.
The analytic realization requires a fiberwise Remez/valency condition in
normalized Fourier coordinates; finite-branch jet surjectivity alone is not
used as a sublevel theorem.  Circular families retain the quantitative jet
condition.  Under the primitive hyperbolicity, single-source,
finite-thickening, branch-section, geometric-flux, split-standard-family,
connected-pullback, and pathwise inputs listed below, the complete first
defect satisfies CM2.  No unconditional generic circular-table claim is made.

\end{abstract}

\section{The exact target}

Throughout, $\{T_s:|s|<s_0\}$ is a $C^3$ compact family of finite-horizon
planar dispersing billiard collision maps generated by a $C^3$ normal
deformation of strictly convex scatterer boundaries and transported to one
reference collision space $M$.  The chosen analytic-boundary or
finite-dimensional circular deformation is subject to
$\mathrm{QT}_{\rm path}$, specified in
Section~\ref{sec:parameter-state}.  The table at $s=0$ contains the
unobstructed clean cycle used in
Proposition~\ref{prop:exact-circular-anosov}.  We consider only the complete
\emph{first} moving-table defect: regular Jacobian/density derivatives and all
same-label physical one-sided faces are assembled before centering.  The tests
$f,g$ belong to $C^r(M)$, with $r$ above the Fourier summability threshold,
and $\mu_0(g)=0$.  Second shape defects and unassembled point atoms are outside
the statement.

Let $P_s$ be the Markov-normalized transfer operator on a fixed collision
fiber, $P_s\one=\one$, let
\[
  \Pi h=\mu_0(h)\one,\qquad Q=P_0-\Pi,
\]
and, for finite $s$, let
\[
  \Pi_s h=\mu_s(h)\one,\qquad Q_s=P_s-\Pi_s,
\]
after transporting the SRB state $\mu_s$ to the reference collision space.
All finite-parameter estimates below are uniform on the fixed compact table
chart.
and let $K_V$ denote the complete first moving-table defect after all regular
terms, same-label physical faces, and grazing blocks have been assembled.  The
two-time estimate needed for a prescribed centered source $g$ is
\begin{equation}
\label{eq:CM2}
 \bigl|\langle Q^nK_VQ^mg,f\rangle\bigr|
 \le C\tau^m\vartheta^n\|g\|_{C^r}\|f\|_{C^r},
 \qquad m,n\ge0,
\end{equation}
for some $\tau,\vartheta<1$.  For the difference quotient it is enough to
have the same estimate on words of depth at most $L(s)$ together with an
$o(1)$ CM2 tail beyond $L(s)$; this is the pathwise diagonal criterion below.

The operator must be centered as a whole:
\begin{equation}
\label{eq:centered-defect}
  K_V=\Pi^\perp\left.\partial_sP_s\right|_{s=0}\Pi^\perp.
\end{equation}
Centering separate faces is neither invariant nor sufficient.  Without
\eqref{eq:centered-defect}, rank-one terms generally give an additive bound
$O(\tau^m+\vartheta^n)$, whose double sum diverges.

\section{Why the direct projective-cone bridge fails}

The projective cone for dispersing billiards controls integrals of a density
against positive dynamically H\"older tests on stable curves, short-leaf
mass, and H\"older comparison of such integrals between nearby stable curves.
A moving inverse-branch boundary instead produces, before physical assembly,
a coefficient of the form
\begin{equation}
\label{eq:face-model}
  h\longmapsto w\,(h\vert_{I^+}-h\vert_{I^-})\,\delta_\Gamma,
\end{equation}
where $I^\pm$ are one-sided terminal preimage arcs.  In the radial collision
atlas these arcs land on the grazing seam and are transversal to the stable
test foliation.  Thus \eqref{eq:face-model} asks for a transverse trace, not a
stable-leaf average.

The following elementary local model rules out the missing implication.

\begin{proposition}[Stable-leaf control does not dominate transverse trace]
\label{prop:no-trace}
Let $R=(-2,2)^2$, let the local stable leaves be the horizontal segments
$W_y=(-2,2)\times\{y\}$, and let the transversal face be
$I=\{0\}\times(-1,1)$.  Fix $0<q<1$ and $0<\gamma<1$.  There are smooth
functions $u_\varepsilon$ supported in an $\varepsilon$-tube about $I$ such
that
\[
 u_\varepsilon\vert_I=b\ne0,
\]
while, uniformly over bounded H\"older test functions $\psi$ on stable
subsegments $W\subset W_y$,
\begin{align}
 \left|\int_Wu_\varepsilon\psi\,dm_W\right|
 &\le C\min\{|W|,\varepsilon\}\|\psi\|_\infty,
       \label{eq:leaf-small}\\
 \sup_{0<|W|\le1}|W|^{-q}
 \left|\int_Wu_\varepsilon\psi\,dm_W\right|
 &\le C\varepsilon^{1-q}\|\psi\|_\infty,\label{eq:short-small}
\end{align}
and the H\"older comparison of these integrals between nearby parallel
stable leaves is $O(\varepsilon)$.  Consequently no norm or order estimate
using only the stable-leaf quantities in
\eqref{eq:leaf-small}--\eqref{eq:short-small} and their transverse H\"older
comparisons can bound $\|u\vert_I\|_{C^0}$.
\end{proposition}

\begin{proof}
Choose $\phi\in C_c^\infty((-1,1))$ with $\phi(0)=1$ and
$b\in C_c^\infty((-1,1))$, $b\ne0$, and set
\[
  u_\varepsilon(x,y)=\phi(x/\varepsilon)b(y).
\]
The trace assertion is immediate.  A horizontal subsegment meets the support
in length at most $C\min\{|W|,\varepsilon\}$, proving
\eqref{eq:leaf-small}.  If $|W|\le\varepsilon$, multiplication by
$|W|^{-q}$ gives at most $C\varepsilon^{1-q}$; if
$|W|\ge\varepsilon$, it gives at most
$C\varepsilon|W|^{-q}\le C\varepsilon^{1-q}$, proving
\eqref{eq:short-small}.  Replacing $y$ by $y'$ changes the integral by at most
$C\varepsilon|y-y'|\|b\|_{C^1}\|\psi\|_{C^0}$, with the usual additional
H\"older test-comparison term.  The trace is independent of
$\varepsilon$, so the asserted domination is impossible.
\end{proof}

\begin{corollary}[Failure of the proposed positive cone operator]
The local face map \eqref{eq:face-model} is not known to be positive for the
Demers--Liverani cone order, and its boundedness cannot be inferred from
projective contraction.  In particular, the implication
\[
 \Theta(P_0^mf_+,P_0^mf_-)\le C\chi^m
 \quad\Longrightarrow\quad
 \|K_VP_0^m(f_+-f_-)\|_Y\le C'\chi^m
\]
is invalid without an additional transverse-trace theorem.
\end{corollary}

\begin{proof}
The cone order is tested through stable-leaf integrals.  The face coefficient
contains precisely the trace excluded by Proposition~\ref{prop:no-trace}.
Projective contraction therefore supplies no order comparison for the two
output face measures.
\end{proof}

This is a go/no-go conclusion, not merely a missing estimate: the direct
``past cone $\to$ face $\to$ future family'' proof must be abandoned unless
physical assembly first converts the terminal trace into an integral
functional controlled by the cone.  The existing adjoint dyadic bound does
not perform that conversion.

\section{A sufficient two-sided bridge}

The complete defect is a kernel located at the event time.  It must be
controlled from the past and the future simultaneously.  The following
criterion is an exact sufficient condition, but not a characterization of
CM2.

\begin{definition}[Centered nuclear recovery bridge]
Let $X_-$ be a Banach space of centered past observables and let $Y_+$ be a
Banach space of centered signed measured unstable families.  A centered
operator $K:X_-\to Y_+$ has a nuclear recovery presentation if
\begin{equation}
\label{eq:nuclear}
 K=\sum_{j\ge1}\ell_j\otimes\eta_j,
 \qquad
 \sum_{j\ge1}\|\ell_j\|_{X_-^*}\|\eta_j\|_{Y_+}<\infty,
\end{equation}
where $\ell_j(\one)=0$ and $\eta_j(\one)=0$ after the complete regular and
physical-face assembly.
\end{definition}

\begin{theorem}[Two-sided bridge implies CM2]
\label{thm:bridge}
Assume
\begin{align}
 \|Q^mg\|_{X_-}&\le C_-\tau^m\|g\|_{X_-},
       \label{eq:past}\\
 |\eta(\cU_0^nf)|&\le C_+\vartheta^n
       \|\eta\|_{Y_+}\|f\|_{C^1},\label{eq:future}
\end{align}
for centered $g$ and centered $\eta$.  If $K$ has the presentation
\eqref{eq:nuclear}, then
\[
 |(KQ^mg)(\cU_0^nf)|
 \le C_-C_+\mathfrak N(K)\tau^m\vartheta^n
       \|g\|_{X_-}\|f\|_{C^1},
\]
where
$\mathfrak N(K)=\sum_j\|\ell_j\|_{X_-^*}\|\eta_j\|_{Y_+}$.
\end{theorem}

\begin{proof}
Using \eqref{eq:nuclear}, \eqref{eq:past}, and \eqref{eq:future},
\begin{align*}
 |(KQ^mg)(\cU_0^nf)|
 &\le\sum_j|\ell_j(Q^mg)|\,
              |\eta_j(\cU_0^nf)|\\
 &\le C_-C_+\tau^m\vartheta^n\|g\|_{X_-}\|f\|_{C^1}
       \sum_j\|\ell_j\|_{X_-^*}\|\eta_j\|_{Y_+}.
\end{align*}
\end{proof}

\begin{remark}
Used globally, a one-gap functional-correlation estimate gives only
$O(\tau^m+\vartheta^n)$ and does not by itself prove CM2.  It is nevertheless
exactly sufficient in the central band, where $m\asymp n$: after smoothing
the event kernel, one decoupling gives $e^{-c\max\{m,n\}}$, hence an
$e^{-c'(m+n)}$ bound there.  Forward and reverse standard-family coupling
handle the two off-diagonal regions; Lemma~\ref{lem:three-region} assembles
the product estimate.  Thus no operator-ideal estimate is used in the
principal proof.
\end{remark}

\begin{remark}[The nuclear criterion is generally too strong]
\label{rem:nuclear-too-strong}
A physical face current is supported on a graph in a stable--unstable product
chart.  Its kernel is therefore analogous to a pullback or a Fourier integral
operator, not to a smoothing operator.  Even after double centering, a graph
delta need not be nuclear between natural past and future spaces.  Thus failure
of \eqref{eq:nuclear} or of the projective-tensor condition below does
\emph{not} imply failure of CM2.  Product decay may instead arise dynamically,
from oscillation and transversality after the graph is acted on by both time
directions.
\end{remark}

\section{Recovery-time weighted future space}

For a measured unstable component $(W,\rho)$ let $R(W)$ be the first time at
which its canonical forward subdivision is a proper standard family.  For
$\omega>0$, define schematically
\begin{equation}
\label{eq:Yplus}
 \|\eta\|_{Y_+}
 :=\inf\sum_j|a_j|e^{\omega R(W_j)}
       \bigl(1+\mathrm{Reg}(\rho_j)\bigr),
 \qquad
 \eta=\sum_ja_j(W_j,\rho_j).
\end{equation}
The infimum also allows the positive and negative Jordan parts to be split by
recovery rank.  Standard-family coupling and recovery give
\eqref{eq:future} once the constants in \eqref{eq:Yplus} are chosen within
the recovery tail.

For a radial grazing block $b$, the physical flux factorization gives a mass
bound $M_b\le C2^{-b}$.  If its geometric deterioration is polynomial in
$2^b$, the recovery lemmas give
\[
 R_b\le C_0+C_1b.
\]
Hence
\begin{equation}
\label{eq:recovery-sum}
 \sum_{b\ge0}M_be^{\omega R_b}<\infty
 \quad\text{whenever}\quad \omega C_1<\log2.
\end{equation}
Thus initial finite $Z$ is unnecessary.  What remains is not future recovery,
but the past functional $\ell_j$ in \eqref{eq:nuclear}.

\subsection{Recovery windows and the actual summability target}

For differentiability of the susceptibility, a uniform pointwise CM2 bound
for every grazing block is stronger than necessary.  The required conclusion
is absolute summability of the complete block series.  This distinction
removes an artificial exponential loss in the grazing rank.

\begin{proposition}[Recovery-window summability]
\label{prop:recovery-window-summability}
Let $B_b(m,n)$ be the scalar pairing produced by block $b$ at past and future
times $m,n$.  Suppose that there are masses $M_b$, recovery times
$R_b^-,R_b^+$, constants $C<\infty$, and $0<\tau,\vartheta<1$ such that
\begin{equation}
\label{eq:shifted-block-decay}
 |B_b(m,n)|
 \le CM_b\tau^{(m-R_b^-)_{+}}
             \vartheta^{(n-R_b^+)_{+}},
 \qquad m,n\ge0.
\end{equation}
Then
\begin{equation}
\label{eq:recovery-window-sum}
 \sum_{b,m,n}|B_b(m,n)|
 \le C'
 \sum_bM_b(1+R_b^-)(1+R_b^+).
\end{equation}
In particular, if $M_b\le C2^{-\alpha b}$ for some $\alpha>0$ and
$R_b^\pm\le C_0+C_1b$, then the full threefold series converges.  The same
conclusion holds uniformly for parameter difference quotients if the
constants and block masses are uniform.

The deterministic recovery times may be replaced by a countable refinement
$B_b=\sum_jB_{b,j}$ with masses $M_{b,j}$ and times $R_{b,j}^\pm$, provided
\begin{equation}
\label{eq:bi-recovery-moment}
 \sum_jM_{b,j}(1+R_{b,j}^-)(1+R_{b,j}^+)
 \le CM_b(1+b)^2.
\end{equation}
It is enough for the two-sided recovery tail, after its deterministic
$O(b)$ entrance time, to be exponential.
\end{proposition}

\begin{proof}
For every $R\ge0$ and $0<q<1$,
\[
 \sum_{k\ge0}q^{(k-R)_+}
 \le R+1+\frac1{1-q}.
\]
Sum \eqref{eq:shifted-block-decay} first in $m,n$ and then in $b$.
Exponential block mass dominates the quadratic recovery window.  Applying the
same calculation to every refined piece gives
\eqref{eq:bi-recovery-moment}.  An exponential tail for
$R_{b,j}^-+R_{b,j}^+-O(b)$ has uniformly bounded polynomial moments, proving
the last assertion.
\end{proof}

\begin{lemma}[Two marginal growth tails imply bi-recovery]
\label{lem:marginal-to-bi-recovery}
Let a block of total variation $M_b$ be cut by the common refinement of its
past and future recovery trees, with refined masses $M_{b,j}$ and recovery
times $R_{b,j}^\pm$.  Assume cutting preserves total variation and, for some
$r_b\le C_0+C_1b$, $0<\rho<1$,
\[
 \sum_{j:R_{b,j}^\pm>r_b+t}M_{b,j}
 \le CM_b\rho^t,
 \qquad t\ge0,
\]
for each sign.  Then \eqref{eq:bi-recovery-moment} holds.
\end{lemma}

\begin{proof}
If $R_{b,j}^-+R_{b,j}^+>2r_b+t$, at least one marginal time is larger than
$r_b+t/2$.  The union bound therefore gives an exponential tail for the sum
of the two recovery times under the normalized weights $M_{b,j}/M_b$.
Its second moment is $O((1+r_b)^2)$.  Since
$(1+R^-)(1+R^+)\le(1+R^-+R^+)^2$, the claim follows.
\end{proof}

For a general graph current the relevant geometric input is therefore a
\emph{proper bi-recovery} statement, not an estimate of a Fourier constant on
the original grazing curve.  For the clean first radial current this input is
Corollary~\ref{cor:radial-bi-recovery}.  A central two-time theorem then gives
\eqref{eq:shifted-block-decay}.  The fully unrecovered rectangle is handled by
total variation, the two mixed rectangles by the corresponding recovered
one-sided estimate, and the fully recovered rectangle by the product bound.
The growth-lemma tail sums these windows through
\eqref{eq:bi-recovery-moment}.

Consequently even the generic $M_b=O(2^{-b})$ assembly margin is sufficient
once proper bi-recovery is proved; the stronger first-radial
$O(2^{-(3-\epsilon_*)b})$ margin gives additional room but is no longer
needed to absorb an exponential Fourier constant in $b$.  This decouples the
grazing rank from the return-rank truncation in
Proposition~\ref{prop:countable-truncation}.

\section{Billiard-specific block problem}

For each genuine physical block $b$, the complete difference quotient must be
written, after artificial-face cancellation, as
\begin{equation}
\label{eq:block-kernel}
 K_{V,b}=\sum_{k\ge1}\ell_{b,k}\otimes\eta_{b,k},
\end{equation}
with
\begin{equation}
\label{eq:block-sum}
 \sum_{b,k}\|\ell_{b,k}\|_{X_-^*}
             \|\eta_{b,k}\|_{Y_+}<\infty.
\end{equation}
The desired estimate follows from Theorem~\ref{thm:bridge}.  Three sufficient
mechanisms for proving \eqref{eq:block-sum} are:
\begin{enumerate}[label=\textup{(\roman*)}]
\item a reverse-time conditional-mixing theorem for the terminal slice,
      furnishing $X_-^*$ directly;
\item a two-sided Young-tower lift of the complete centered defect, with an
      event kernel that is nuclear between the past and future tower norms;
\item a smoothing/nuclear decomposition of the terminal restriction, with a
      quantitative bunching inequality strong enough to sum the smoothing
      error over the homogeneity blocks.
\end{enumerate}
The projective-cone contraction alone supplies none of these three inputs.
These mechanisms are not exhaustive because a non-nuclear graph current can
still equidistribute under the two-sided dynamics.

\subsection{A two-sided tower formulation}

The natural extension gives a sharper version of alternative (ii).  Let
$\Delta$ be a two-sided Young tower for the fixed billiard, with past and
future cylinder filtrations $\mathcal F_i^-$ and $\mathcal F_j^+$.  Denote
the associated martingale differences by
\[
 D_i^-=\mathbb E(\,\cdot\mid\mathcal F_i^-)
       -\mathbb E(\,\cdot\mid\mathcal F_{i-1}^-),
 \qquad
 D_j^+=\mathbb E(\,\cdot\mid\mathcal F_j^+)
       -\mathbb E(\,\cdot\mid\mathcal F_{j-1}^+).
\]
A lifted event kernel $\widehat K$ has summable mixed cylinder variation if
\begin{equation}
\label{eq:mixed-cylinder}
 \sum_{i,j\ge1}
 \|D_i^-D_j^+\widehat K\|_{\mathcal B_-^*\widehat\otimes_\pi
                                      \mathcal B_+}<\infty.
\end{equation}
The indices start at one: complete double centering removes the two axis
terms $D_i^-D_0^+\widehat K$ and $D_0^-D_j^+\widehat K$.

\begin{proposition}[Mixed-cylinder criterion]
\label{prop:mixed-cylinder}
Suppose the two-sided tower has exponential tails and the past and future
transfer actions contract the centered martingale blocks at rates
$\tau^m$ and $\vartheta^n$.  If the complete lifted moving defect satisfies
\eqref{eq:mixed-cylinder}, then it has a nuclear recovery presentation
\eqref{eq:nuclear} and hence satisfies CM2.
\end{proposition}

\begin{proof}
Expand
\[
 \widehat K=\sum_{i,j\ge1}D_i^-D_j^+\widehat K
\]
in the projective tensor product.  Each summand is a norm-convergent sum of
rank-one past/future tensors by the definition of
$\widehat\otimes_\pi$.  Summability in \eqref{eq:mixed-cylinder} therefore
gives \eqref{eq:nuclear}.  Shifting $m$ steps into the past and $n$ steps into
the future contracts the two centered tensor factors separately.  Apply
Theorem~\ref{thm:bridge}.
\end{proof}

This criterion identifies one strong geometric issue.  A smooth two-dimensional
density on a tower rectangle has summable mixed cylinder variation.  A face
current is instead supported on a one-dimensional graph in the local
stable--unstable product chart.  Its graph delta is not automatically in the
projective tensor product.  Proving \eqref{eq:mixed-cylinder} is therefore a
genuine conditional-mixing theorem, not routine tower bookkeeping.  The
finite-strip central face supplies uniform transversality and one-sided
recovery, but these facts alone do not prove the mixed tensor estimate.

\section{The intrinsic dynamic graph-current formulation}

Let $\gamma_b:I_b\to M$ parametrize one side of a complete physical face in
block $b$, let $F_b\gamma_b$ denote its assembled terminal image, and let
$w_b$ be the signed coarea weight after artificial-face cancellation.  The
face contribution to CM2 has the model form
\begin{equation}
\label{eq:dynamic-graph}
 B_{m,n}^{(b)}(g,f)
 =\int_{I_b} g(T^{-m}\gamma_b(t))
                  f(T^nF_b\gamma_b(t))w_b(t)\,dt .
\end{equation}
Equivalently, if $J_{V,b}$ is the graph current on $M\times M$ induced by
$t\mapsto(\gamma_b(t),F_b\gamma_b(t))$, then
\begin{equation}
\label{eq:product-current}
 B_{m,n}^{(b)}(g,f)
 =\bigl((T^{-m}\times T^n)_*J_{V,b}\bigr)(g\otimes f).
\end{equation}
Complete double centering removes the two marginals of $J_{V,b}$.  The
intrinsic CM2 statement is therefore two-parameter equidistribution of a
centered graph current under $T^{-1}\times T$ at a product rate.

This formulation explains both the difficulty and the opportunity.  A graph
curve is one-dimensional, whereas the unstable bundle of the product map is
two-dimensional, so ordinary standard-family coupling does not apply
directly.  On the other hand, the current need not be statically nuclear:
hyperbolic stretching sends its Fourier frequencies to exponentially large
scales, and polynomial Fourier decay of the induced slice is already enough
to produce exponential decay in collision time.

\begin{proposition}[Fourier reduction on one induced rectangle]
\label{prop:fourier-reduction}
Assume that an induced return map has a stable--unstable product chart in
which the centered face slice induces a finite signed measure $\nu_b$ on the
contracting coordinate.  Suppose that, uniformly over admissible inverse
words and over the relevant $C^2$ phase maps $F$,
\begin{equation}
\label{eq:fourier-decay}
 |\widehat{F_*\nu_b}(\xi)|\le C_b(1+|\xi|)^{-\eta}
\end{equation}
for some $\eta>0$.  Suppose also that, after excluding phase cancellation,
the expanding frequencies generated by $m$ past returns and $n$ future
returns are bounded below by
\[
 c\max\{\Lambda_-^m,\Lambda_+^n\},
\]
apart from cylinder pieces whose total centered mass is bounded by
$C\tau_0^m\vartheta_0^n$.  Then, for sufficiently regular
centered cylinder observables,
\begin{equation}
 |B_{m,n}^{(b)}(g,f)|
 \le C C_b\,\tau^m\vartheta^n\|g\|_{C^r}\|f\|_{C^r}
\end{equation}
for some $\tau,\vartheta<1$.
\end{proposition}

\begin{proof}
Decompose $g$ and $f$ into Fourier modes in the product chart and split the
sum into low and high frequencies.  Centering and cylinder mixing control the
low-frequency part.  In the high-frequency part, hyperbolicity converts the
two return depths into the lower frequency scale
$c\max\{\Lambda_-^m,\Lambda_+^n\}$.  Since
\[
 \max\{\Lambda_-^m,\Lambda_+^n\}
 \ge (\Lambda_-^m\Lambda_+^n)^{1/2},
\]
\eqref{eq:fourier-decay} gives the product factor
$\Lambda_-^{-\eta m/2}\Lambda_+^{-\eta n/2}$.  Summability of the Fourier
coefficients of $C^r$ tests and the exceptional-cylinder bound complete the
estimate.  The loss of one half in each exponent is harmless.
\end{proof}

This preliminary proposition is not used as an assumption in the final proof.
Uniform \eqref{eq:fourier-decay} is supplied below by the cohomological
local slope-energy theorem, the dominant-coordinate reduction, and the
physical fractional nonstationary-phase estimate.  One-sided mixing alone
would not supply it.

\subsection{A three-region reduction}

Fourier analysis is not needed uniformly over the whole $(m,n)$ quadrant.
One-sided standard-family recovery gives a forward estimate with a past
regularity cost, and reversibility gives its backward analogue.  The following
elementary lemma isolates the remaining diagonal band.

\begin{lemma}[Off-diagonal coupling plus diagonal oscillation]
\label{lem:three-region}
Suppose
\begin{align}
 |B_{m,n}|&\le C e^{\beta_-m-\alpha_+n},\label{eq:forward-cone}\\
 |B_{m,n}|&\le C e^{-\alpha_-m+\beta_+n},\label{eq:backward-cone}
\end{align}
with all four exponents positive.  Choose
$A>\beta_-/\alpha_+$ and $D>\beta_+/\alpha_-$.  Suppose in addition that
\[
 |B_{m,n}|\le Ce^{-\gamma\max\{m,n\}}
\]
whenever $D^{-1}m<n<Am$.  Then there are
$\tau,\vartheta\in(0,1)$ such that
$|B_{m,n}|\le C'\tau^m\vartheta^n$ for all $m,n$.
\end{lemma}

\begin{proof}
If $n\ge Am$, \eqref{eq:forward-cone} is bounded by
$e^{-c_+(m+n)}$ for some $c_+>0$.  If $m\ge Dn$, use
\eqref{eq:backward-cone}.  The remaining indices lie in the stated diagonal
band, where $\max\{m,n\}\ge(m+n)/2$.  Taking the minimum of the three positive
exponents proves the claim.
\end{proof}

For billiards, \eqref{eq:forward-cone} follows from signed
standard-family coupling after cutting the past-weighted density into
homogeneous pieces: the growth lemma gives the coupling factor
$e^{-\alpha_+n}$, while differentiating the $m$-step past weight costs at
most $e^{\beta_-m}$.  The marked-amplitude moment sums the homogeneous
pieces.  Applying the same argument to the reversible collision map gives
\eqref{eq:backward-cone}.  All constants are uniform on the fixed compact
atlas; recovered grazing pieces enter only after their recorded recovery
times; these are the signed standard-family forms of the usual growth and
coupling estimates \cite{Chernov1999,DLreview}.
Consequently the genuinely new slope/Fourier theorem is needed only for
$m\asymp n$.

This also corrects the naive ambient phase test.  The graph
$t\mapsto(T^{-m}\gamma(t),T^nF\gamma(t))$ is one-dimensional in a
four-dimensional product.  Its first two derivatives cannot be uniformly
transverse to every Fourier covector.  In the diagonal band one must first
discard the contracted coordinates and pass to the dominant past-stable and
future-unstable coordinates.  The resulting two-dimensional plane curve is
governed by the stable--unstable temporal-distance/holonomy function.  A
finite-difference slope-value QNL is appropriate for this reduced phase.  If
one insists on working in the full four-dimensional chart, a finite-type
condition involving derivatives through order at least four is required.

\subsection{A verified radial nonlinearity seed}

The first radial face is not an abstract graph.  For two circular scatterers
with radii $R_j,R_i$ and centre distance $d>R_j+R_i$, put the source centre at
the origin, the target centre at $(d,0)$, and parametrize the source boundary
by
\[
 q_j(\alpha)=R_j(\cos\alpha,\sin\alpha).
\]
Let $r(\alpha)=|(d,0)-q_j(\alpha)|$ and
$\beta(\alpha)=\arg((d,0)-q_j(\alpha))$.  The two directions tangent to the
target circle are
\[
 \theta_\sigma(\alpha)
 =\beta(\alpha)+\sigma\arcsin\frac{R_i}{r(\alpha)},
 \qquad \sigma\in\{-1,1\}.
\]
Thus the radial physical face in Birkhoff coordinates is locally
$\varphi_\sigma(\alpha)=\theta_\sigma(\alpha)-\alpha$.

\begin{lemma}[The circular tangency face is genuinely curved]
\label{lem:circular-face-curvature}
At the collinear source point $\alpha=0$, write $r_0=d-R_j$.  Then
\begin{equation}
\label{eq:circular-face-curvature}
 \varphi_\sigma''(0)
 =-\sigma\frac{R_i dR_j}
 {r_0^2\sqrt{r_0^2-R_i^2}}\ne0.
\end{equation}
Moreover the angular coarea derivative of
\[
 G=|((d,0)-q_j)\times v_\theta|^2-R_i^2
\]
is nonzero there.  Hence each tangency choice gives a smooth, non-affine
central face germ with a smooth coarea density.
\end{lemma}

\begin{proof}
Direct differentiation gives
\[
 r'(0)=0,\qquad r''(0)=\frac{dR_j}{r_0},\qquad
 \beta''(0)=0.
\]
Since
\[
 \frac d{dr}\arcsin\frac{R_i}{r}
 =-\frac{R_i}{r\sqrt{r^2-R_i^2}},
\]
equation \eqref{eq:circular-face-curvature} follows.  On the tangency face,
$|(d,0)-q_j|\times v_\theta|=R_i$ while the longitudinal component has
magnitude $\sqrt{r_0^2-R_i^2}$; therefore
$|\partial_\theta G|=2R_i\sqrt{r_0^2-R_i^2}>0$.
\end{proof}

This proves a useful but limited fact: the radial face supplies a concrete
one-step nonlinearity seed and no Baker-type affine resonance is forced by
the circular geometry.  It does not prove uniform quantitative nonlinearity
after arbitrary past and future words.  Near the diagonal where the two
terms in the reduced phase
\begin{equation}
\label{eq:phase}
 \Phi_{m,n,\xi,\zeta}(t)
 =\xi\cdot T^{-m}\gamma(t)+\zeta\cdot T^nF\gamma(t)
\end{equation}
have comparable size, their derivatives can cancel;
this is the only time regime in which a genuine Dolgopyat/sum-product
argument is needed.  Off that diagonal one time direction dominates and
ordinary nonstationary phase applies.

\subsection{A countable-alphabet truncation lemma}

The latest smooth result closest to the required mechanism is Leclerc's
Fourier-decay theorem for equilibrium states of nonlinear, area-preserving
smooth Axiom--A surface diffeomorphisms.  Its quantitative nonlinearity
condition is expressed through a temporal-distance function and is generic,
with an explicit periodic-point cocycle test.  A billiard return scheme is
piecewise smooth with a countable alphabet, so that theorem is not directly
applicable.  The countable alphabet itself, however, is not the fundamental
obstruction.

\begin{proposition}[Exponential-tail reduction to uniform finite truncations]
\label{prop:countable-truncation}
Let an induced symbolic system have a return/singularity rank $R$ with
\[
 \mu\{R>L\}\le C e^{-aL}.
\]
Suppose a Fourier proof at frequency $|\xi|$ uses at most
$N(\xi)\le c_0\log(2+|\xi|)$ successive symbols.  Assume that the subsystem
with all ranks at most $L$ satisfies, uniformly in the truncation,
\[
 \left|\int e^{i\xi\psi}\chi\,d\mu_L\right|
 \le C e^{cL}|\xi|^{-\eta}.
\]
Then the full countable system has polynomial Fourier decay, provided
$a,\eta>0$ and the finite-truncation constants have the displayed
exponential bound.
\end{proposition}

\begin{proof}
The mass of orbit words which see a rank larger than $L$ is at most
$C N(\xi)e^{-aL}$.  On the complement apply the finite-truncation estimate.
Choose $L=\delta\log|\xi|$ with $0<\delta<\eta/c$ (with no upper restriction
if $c=0$).  The good contribution is
$O(|\xi|^{-\eta+c\delta})$ and the bad contribution is
$O(\log|\xi|\,|\xi|^{-a\delta})$.  Both have a positive power saving.
\end{proof}

Thus this fallback would require a \emph{uniform} singular version of
Leclerc's quantitative-nonlinearity estimate on finite homogeneity/return
truncations.  Exponential tower tails could then remove the truncation, but
the truncation-uniform constants are not supplied by that route.  The
direct gate-resolved slope argument below bypasses this fallback entirely.

\subsection{Uniform QNL requires anchored big images}

Inspection of the QNL proof shows that a single nonzero periodic cocycle is
not, by itself, a uniform truncation theorem.  The power exponent also depends
on a cutting constant: at every scale one removes a child interval carrying
at least an $\eta_{\rm cut}$ fraction of its parent's conditional mass.  For
an arbitrary exhaustion by finite subshifts, $\eta_{\rm cut}$ may tend to
zero; then the QNL exponent tends to zero and
Proposition~\ref{prop:countable-truncation} is useless.

The billiard inducing scheme has a possible cure.  A Gibbs--Markov return map
has a finite big-images/preimages core, and on a compact central magnet the
SRB conditional densities on homogeneous unstable leaves are bounded above
and below.  The truncations must therefore be \emph{anchored}: retain a fixed
finite family of full-image connector branches, including the nonlinear
periodic witness, and add all branches up to rank $L$.  The desired uniform
cutting statement is
\begin{equation}
\label{eq:uniform-child-mass}
 \nu_{x,L}(I_{\mathrm{child}})
 \ge\eta_0\,\nu_{x,L}(I_{\mathrm{parent}}),
 \qquad \eta_0>0
\end{equation}
with $\eta_0$ independent of $x,L$, for at least one QNL-separated child at
each scale.  If \eqref{eq:uniform-child-mass} and a uniform UNI constant
$\kappa_0>0$ hold, Leclerc's tree argument gives a QNL exponent
$\Gamma>0$ independent of $L$; only the multiplicative constant is allowed
to grow like $e^{cL}$.

R\"uhr's pressure inequalities quantify weak convergence of Gibbs measures
on countable Markov shifts to equilibrium measures on finite subsystems by
the square root of the pressure gap.  They support the truncation step, but
do not prove \eqref{eq:uniform-child-mass} or uniform UNI.  Strong positive
recurrence/BIP controls the symbolic escape; the billiard geometry must still
provide the nonlinear separated child.

For the collision SRB measure on a compact homogeneous magnet, the unstable
conditional is better than a general fractal Gibbs conditional: in arclength
coordinates its density is bounded above and below.  In that case the
measure-theoretic part of the QNL argument is elementary.

\begin{lemma}[Uniform interval UNI implies power non-concentration]
\label{lem:interval-uni-qnl}
Let $\nu_L$ be measures on an interval whose densities satisfy
$0<c_-\le d\nu_L/dt\le c_+<\infty$, uniformly in $L$.  Let $X_L$ be a family
of functions with the following scale-invariant UNI property: for every
interval $J$ of length $\sigma<\sigma_0$ and every polynomial $P$ in a fixed
finite-dimensional class, there is a subinterval
$J'\subset J$ with $|J'|\ge\kappa|J|$ such that
\[
 |X_L(t)-P(t)|\ge\kappa|J|\quad(t\in J'),
\]
where $\kappa>0$ is independent of $J,P,L$.  Then for every fixed
$\alpha>0$ there are $C,\Gamma>0$, independent of $L$, such that
\[
 \nu_L\{t\in J:|X_L(t)-P(t)|\le |J|^{1+\alpha}\}
 \le C|J|^\Gamma\nu_L(J).
\]
\end{lemma}

\begin{proof}
Partition every surviving interval into a bounded number of comparable
children and discard a child contained in the UNI interval $J'$.  The density
bounds imply that at least a fixed fraction $\eta_0>0$ of the conditional
mass is discarded at each generation.  Iterate until the child scale is
comparable to $|J|^{1+\alpha}$.  The number of generations is
$c_\alpha|\log|J||+O(1)$, so the surviving mass is at most
$(1-\eta_0)^{c_\alpha|\log|J||}\nu_L(J)
\le C|J|^\Gamma\nu_L(J)$.
\end{proof}

Thus for the central homogeneous face the principal task can be reduced
further: establish a uniform geometric UNI statement for the
stable--unstable temporal-distance function.  Finite-subsystem approximation
is then needed for long singular itineraries and difference quotients, but
not to manufacture the QNL exponent on the central magnet.

\subsection{A fixed two-branch UNI certificate}

Bertozzi's 2026 countable-Markov exponential-mixing construction provides a
useful template: after sufficient inducing, two fixed inverse branches are
repeated and the derivative of their Birkhoff-roof difference is bounded away
from zero uniformly in the iteration depth.  The same device can be formulated
for the billiard temporal-distance increment.

Let $h_0,h_1$ be two regular inverse return branches on the central magnet,
with a common contraction bound $\lambda<1$, and let $\tau$ denote the
reduced stable--unstable phase increment (for the geometric application,
$\tau=\log J^uF$ after the stable coboundary normalization).  Define
\begin{equation}
\label{eq:two-branch-temporal-distance}
 \Psi_N(x)=\sum_{k=1}^N
 \bigl[\tau(h_0^kx)-\tau(h_1^kx)\bigr],
 \qquad
 \Psi_\infty=\lim_{N\to\infty}\Psi_N .
\end{equation}

\begin{proposition}[Anchored two-branch certificate]
\label{prop:two-branch-uni}
Suppose the series in \eqref{eq:two-branch-temporal-distance} converges in
$C^1$ and
\[
 \inf_{x\in I}|\Psi_\infty'(x)|\ge c_0>0
\]
on a central unstable interval $I$.  Then, for all sufficiently large $N$,
\[
 \inf_{x\in I}|\Psi_N'(x)|\ge c_0/2.
\]
If the two branches and their fixed connector branches are retained in every
anchored truncation, this UNI constant is independent of the truncation rank.
\end{proposition}

\begin{proof}
$C^1$ convergence makes
$\|\Psi_N'-\Psi_\infty'\|_\infty<c_0/2$ for all large $N$.  The second
assertion follows because the calculation uses only the anchored branches,
not the additional symbols admitted by a larger truncation.
\end{proof}

\begin{corollary}[First-term domination after acceleration]
\label{cor:first-term-uni}
Assume $\|Dh_i\|\le\lambda<1$, $\|D\tau\|_\infty\le M$, and
\[
 \inf_I|D(\tau\circ h_0-\tau\circ h_1)|\ge d_0\lambda .
\]
If
\[
 \frac{2M\lambda^2}{1-\lambda}\le\frac12d_0\lambda,
\]
then $\inf_I|\Psi_\infty'|\ge d_0\lambda/2$.
\end{corollary}

\begin{proof}
The derivative of the sum of all terms with $k\ge2$ is bounded by
$2M\sum_{k\ge2}\lambda^k=2M\lambda^2/(1-\lambda)$.  Subtract this from the
first-term lower bound.
\end{proof}

This changes the next billiard calculation from an infinite statistical
problem into a finite geometric one.  Choose two central return words whose
one-step radial tangency germs have different curvature coefficients in
\eqref{eq:circular-face-curvature}; attach fixed regular connectors to the
magnet; and compute $\Psi_\infty'$ (or a rigorously bounded finite
approximation plus its geometric tail).  A nonzero lower bound is open under
small table perturbations and gives the uniform UNI input for
Lemma~\ref{lem:interval-uni-qnl}.  The non-equality of the two one-step
curvatures is a promising seed, but it is not by itself the required lower
bound because the stable coboundary and all subsequent return terms enter
$\Psi_\infty$.
After accelerating both branches, Corollary~\ref{cor:first-term-uni} shows
that it is enough to retain a lower bound of order $\lambda$ for the
first-term curvature separation; the remaining infinite temporal-distance
tail is only $O(\lambda^2)$.  This remains a useful branchwise fallback.  The
next calculation gives a more direct fixed periodic witness for the smooth
finite truncations.

\subsection{An exact Anosov-cocycle witness from a normal two-cycle}

There is a more direct way to certify the geometric nonlinearity needed by
the smooth finite truncations.  Leclerc's criterion for a smooth
area-preserving Axiom--A surface map says that QNL follows if the stable or
unstable distribution is not $C^2$.  At a hyperbolic fixed point, this can be
checked by the Anosov cocycle.  In determinant-one adapted coordinates, write
the map as $(F,G)$ with derivative
$\operatorname{diag}(\lambda,\mu)$, where $\lambda>1$, $\mu=\lambda^{-1}$.
The relevant scalar is
\begin{equation}
\label{eq:anosov-cocycle}
 \mathfrak A(F,G)
 =\frac{G_{xyy}}{\mu}
 -\frac{F_{xy}G_{xy}}{\lambda-1}
 -\frac{G_{xx}F_{yy}}{1-\mu^3}
 -\frac{G_{xy}G_{yy}}{\mu^2},
\end{equation}
with all derivatives evaluated at the fixed point.

\begin{proposition}[Exact circular normal-cycle witness]
\label{prop:exact-circular-anosov}
Consider two unit circular scatterers whose facing boundary points are
separated by a unit free-flight gap.  Let $\mathcal R=T^2$ be the billiard
return map to the first component near the unobstructed normal period-two
orbit, in Birkhoff coordinates $(s,p)$ with $p=\sin\varphi$.  Then
\[
 D\mathcal R(0)=
 \begin{pmatrix}7&4\\12&7\end{pmatrix}.
\]
For the determinant-one eigenbasis
\[
 S=\begin{pmatrix}
       1&-\sqrt3/6\\[2pt]
       \sqrt3&1/2
    \end{pmatrix},
 \qquad \det S=1,
\]
write $S^{-1}\mathcal R S=(F,G)$.  Then
\begin{align*}
 \lambda&=7+4\sqrt3,&\mu&=7-4\sqrt3,\\
 F_{xy}=F_{yy}=G_{xx}=G_{xy}=G_{yy}&=0,
 &G_{xyy}&=-\mu.
\end{align*}
Consequently
\[
 \mathfrak A(F,G)=-1.
\]
In particular, the local return dynamics has a nonzero QNL witness which is
stable under sufficiently small clean perturbations of the table.
\end{proposition}

\begin{proof}
Place the circle centres at $(0,0)$ and $(3,0)$ and parametrize the first
boundary near the facing point by $q(s)=(\cos s,\sin s)$.  For outgoing
Birkhoff data $(s,p)$, set
\[
 n=(\cos s,\sin s),\qquad t=(-\sin s,\cos s),\qquad
 v=\sqrt{1-p^2}\,n+pt.
\]
For a ray based at $q$ and a target circle with centre $c$ and radius one,
put $w=q-c$, $b=w\cdot v$, and $c_0=|w|^2-1$.  The first collision time and
reflected velocity are
\[
 t_*=-b-\sqrt{b^2-c_0},\qquad
 v^+=v-2(v\cdot n_*)n_*,
 \quad n_*=q+t_*v-c.
\]
Apply this formula first to the second circle and then to the first.  Taylor
expansion through degree three gives the displayed linear matrix.  Reflection
in the line of centres conjugates $(s,p)$ to $(-s,-p)$ and commutes with the
two-collision return, so the return map is odd and all its quadratic jets
vanish.  Conjugating the cubic jet by $S$ gives
$G_{xyy}=-(7-4\sqrt3)$; substitution in
\eqref{eq:anosov-cocycle} gives $-1$.
\end{proof}

The calculation is scale free: if both radii and the free-flight gap are
multiplied by the same positive number, use dimensionless arclength $s/R$ and
obtain the same return jet.  Thus the local configuration can be inserted at
an arbitrarily small geometric scale into a finite-horizon table, provided
the normal two-cycle remains unobstructed.  More generally, on any fixed clean
finite itinerary the return jet and $\mathfrak A$ depend real-analytically on
the circular radii, centres, and collision points.  Proposition
\ref{prop:exact-circular-anosov} therefore implies that the nonzero-cocycle
tables form an open dense subset of every connected analytic parameter chart
containing this sample.

For an anchored finite horseshoe which contains this normal two-cycle and a
regular connector branch, the same fixed point occurs in every larger
anchored truncation.  Hence the geometric nonlinearity witness itself no
longer deteriorates with truncation rank.  What still requires proof is a
uniform smooth extension/Markov realization of those finite billiard cores,
uniform cutting mass for the SRB conditionals, and transfer of Leclerc's QNL
bound from the equilibrium state to the weighted double-centered face slice.
The exact cocycle closes the local geometry, not these singular-exhaustion
steps.

\subsection{Transfer of QNL to a central weighted face slice}

The equilibrium-state/face-current mismatch is harmless on one compact
homogeneous return rectangle.  Let $W^u$ be a reference unstable leaf and let
$H:\Gamma\to W^u$ be stable holonomy from a central smooth face arc.  Write
$m_\Gamma$ for arclength and let $w$ include the physical face velocity,
coarea factor, and the smooth branch Jacobians.

\begin{lemma}[Weighted face-slice transfer]
\label{lem:weighted-face-qnl-transfer}
Assume, uniformly on a compact clean family, that
\begin{enumerate}[label=\textup{(\roman*)}]
\item $H$ is a $C^{1+\alpha}$ diffeomorphism onto its image and
      $0<c_H\le JH\le C_H$;
\item the SRB conditional $\mu^u$ on $W^u$ has arclength density between
      $c_u>0$ and $C_u<\infty$;
\item $w\in C^\alpha(\Gamma)$ and $\|w\|_\infty\le C_w$.
\end{enumerate}
Define the signed face measure
\[
 \nu=H_*(w\,dm_\Gamma).
\]
If, on every interval $J\subset W^u$, the temporal-distance phase $X$
satisfies
\[
 \mu^u\{t\in J:|X(t)-P(t)|\le\epsilon\}
 \le C_0\epsilon^\Gamma\mu^u(J),
\]
then
\[
 |\nu|\{t\in J:|X(t)-P(t)|\le\epsilon\}
 \le C_1\epsilon^\Gamma |J|,
 \qquad
 C_1=C_0\frac{C_wC_u}{c_Hc_u},
\]
after increasing $C_1$ by a uniform chart constant.  The same conclusion
holds for a finite signed sum of physical face pieces, with $C_1$ multiplied
by the sum of their total-variation bounds.
\end{lemma}

\begin{proof}
The change-of-variables formula gives
\[
 \frac{d|\nu|}{dm_{W^u}}(H z)
 =\frac{|w(z)|}{JH(z)}\le\frac{C_w}{c_H}.
\]
The density bounds for $\mu^u$ compare arclength and $\mu^u$ in both
directions.  Apply the assumed small-value estimate and make those two
comparisons.  Total variation and the triangle inequality handle signed
pieces.
\end{proof}

Lemma~\ref{lem:circular-face-curvature} supplies the central coarea
transversality, while standard clean-rectangle holonomy estimates supply
(i)--(ii).  Thus, once the exact cocycle is placed in an anchored smooth
horseshoe, the central weighted face slice inherits the same QNL exponent.
For pieces whose return word or homogeneity rank tends to infinity, this
local transfer is combined below with the central marked moment and the
two-sided recovery sum; Proposition~\ref{prop:global-labelled-ledger} keeps
those constants on one common refinement.

\subsection{A countable-IFS theorem removes the return truncation}
\label{subsec:countable-ifs}

There is a more direct route than passing through smooth finite horseshoes.
Leclerc--Paukkonen--Sahlsten~\cite{LPS} prove power Fourier decay for Gibbs measures of
finite or countable $C^{1+\alpha}$ conformal IFSs, allowing overlaps.  In the
countable case their assumptions are uniform contraction, bounded distortion,
summable variations, and the light-tail condition
\begin{equation}
\label{eq:lps-light-tail}
 \sum_{a}\sup_{x\in I}
 e^{\phi(h_a x)}|h_a'(x)|^{-\gamma_0}<\infty
 \qquad\text{for some }\gamma_0>0.
\end{equation}
Their Theorem~1.1 says that QNL gives power Fourier decay, including for
H\"older amplitudes and nonstationary phases with a small polynomial loss in
their norms.  Sections~6.1--6.2 prove QNL by finite-difference cutting
arguments for finite MNL and for countable UNI systems with a light tail.  No
finite-alphabet exhaustion or smooth extension of the branch maps is used.

This theorem is unusually well matched to the present induced problem.  Let
$F:I\to I$ be the unstable quotient of a full-branch Gibbs--Markov return to
a central billiard magnet and let $\{h_a:I\to I\}_{a\in\mathcal A}$ be its
inverse branches.  After accelerating once if necessary, the desired
dictionary is
\begin{equation}
\label{eq:lps-dictionary}
 \begin{array}{c|c}
 \text{countable-IFS input}&\text{billiard input}\\ \hline
 h_a&\text{homogeneous inverse return branch}\\
 \phi&\text{normalized induced unstable-Jacobian potential}\\
 \text{uniform contraction}&\text{uniform hyperbolicity after inducing}\\
 \text{bounded distortion}&\text{homogeneous return distortion}\\
 \text{light tail}&\text{return tail plus a small grazing moment}\\
 \text{QNL}&\text{a fixed clean finite-difference gate}.
 \end{array}
\end{equation}

There is a small but important point in the first line of this dictionary.
The one-sided moving-face labels must not be inserted by cutting every return
branch into labelled subbranches: those pieces need not have full image and
would destroy the hypothesis to which the countable-IFS theorem is applied.
The correct object is a full-branch dynamical quotient with a marked amplitude
fiber.

For an interval amplitude use the affine scale-normalized norm
\begin{equation}
\label{eq:scaled-amplitude-norm}
 \|w\|_{\mathcal A^\alpha(J)}
 :=\|w\|_{\infty,J}+|J|^\alpha[w]_{C^\alpha(J)}.
\end{equation}
It is unchanged up to a fixed factor by affine normalization and cannot
increase under restriction.  Every connected component is charged
separately; hence endpoint and component complexity is visible in the sum
rather than hidden in a piecewise-H\"older notation.

\begin{proposition}[Marked full-branch quotient]
\label{prop:marked-full-branch-quotient}
Let $\Lambda$ be a Young return rectangle with product structure and let
$\{\Lambda_a\}_{a\in\mathcal A}$ be a return partition such that
$F|_{\Lambda_a}=T^{R_a}$ maps the stable subset $\Lambda_a$ across $\Lambda$
with the usual stable/unstable Markov property.  Fix a reference unstable
leaf $I$ and let $\pi^s:\Lambda\to I$ be stable holonomy.  Then
\[
 \bar F\circ\pi^s=\pi^s\circ F
\]
defines a one-state full-branch expanding map: for
$I_a=\pi^s(\Lambda_a)$, the map $\bar F:I_a\to I$ is onto and has an inverse
branch $h_a:I\to I_a$.

Suppose now that the lift of a complete physical face through $\Lambda_a$ is
refined into one-sided marked pieces $E_{a,\ell}$, after artificial-face
cancellation.  On the stable quotient its action can be written
\begin{equation}
\label{eq:marked-amplitude-disintegration}
 \nu(\varphi)=
 \sum_{a,\ell}\int_{J_{a,\ell}}
     \varphi(h_a x)w_{a,\ell}(x)\,d\mu^u(x),
\end{equation}
where the $J_{a,\ell}\subset I$ are amplitude pieces and the sign is included
in $w_{a,\ell}$.  Thus the quotient dynamical alphabet remains
$\{h_a\}_{a\in\mathcal A}$; the face label $\ell$ is a fiber mark and does
not have to define a full-image branch.  If
\[
 \sum_{a,\ell}\|w_{a,\ell}\|_{\mathcal A^\alpha(J_{a,\ell})}
      \mu^u(J_{a,\ell})<\infty,
\]
then the marked current is a summable piecewise-H\"older amplitude over the
same full-branch Gibbs system.
\end{proposition}

\begin{proof}
The Markov property says that every unstable leaf in $\Lambda_a$ is mapped
across every stable leaf of $\Lambda$.  Quotienting stable leaves therefore
makes $\bar F|_{I_a}$ onto $I$; contraction of inverse branches and bounded
distortion are the usual quotient estimates.  Disintegrate the coarea face
measure first over return atoms and then over its one-sided physical marks.
Stable holonomy transports each piece to the reference leaf and gives
\eqref{eq:marked-amplitude-disintegration}.  No restriction
$h_a|_{J_{a,\ell}}$ is declared to be a dynamical branch.  The last assertion
is the triangle inequality for the H\"older-amplitude norm.  Matched
artificial cuts contribute opposite weights and disappear before this
disintegration.
\end{proof}

Consequently the standard Young quotient and the moving-face bookkeeping are
compatible provided the marked-amplitude sum is finite.  The construction of
the quotient is standard.  Proposition~\ref{prop:central-marked-moment}
supplies this summability on the central return alphabet, while the
bi-recovery refinement supplies the grazing tail.

The SRB potential part of this audit is now supported by a direct billiard
result.  Lemmas~5.2--5.3 of Lima--Obata--Poletti~\cite{LOP} prove that the
collision Lebesgue/SRB measure is an equilibrium state of the geometric
potential
\[
 \phi=-\log\|DT|_{E^u}\|,
 \qquad P(\phi)=0,
\]
and that $\phi$ becomes H\"older on their countable Markov coding for
finite-horizon dispersing billiards.  Hence the induced potential has
summable variations.  This does not by itself give the one-state return tail,
but it removes the former ambiguity between an MME coding and the required
collision-SRB Gibbs weights.

\begin{proposition}[Central amplitude moment]
\label{prop:central-marked-moment}
Assume the Young return scheme has
\[
 \sum_a\sup_I p_a\,e^{\epsilon R_a}<\infty
\]
and restrict first to a compact regular homogeneity core.  For a complete
first moving-face defect, suppose every collision insertion has uniformly
bounded $C^\alpha$ coarea and stable-holonomy factors.  Then the marked
amplitude disintegration \eqref{eq:marked-amplitude-disintegration} satisfies
\[
 \sum_{a,\ell}e^{\epsilon'R_a}
   \|w_{a,\ell}\|_{\mathcal A^\alpha(J_{a,\ell})}
   \mu^u(J_{a,\ell})<\infty
 \qquad(0<\epsilon'<\epsilon).
\]
After high-homogeneity insertions are assigned to the bi-recovery blocks of
Corollary~\ref{cor:radial-bi-recovery}, the same conclusion holds for the
entire assembled first defect.
\end{proposition}

\begin{proof}
A return word of length $R_a$ has at most $R_a$ possible physical insertion
times.  On the compact regular core each insertion contributes at most a
uniform constant times the Gibbs branch weight.  Hence the left-hand side is
bounded by
\[
 C\sum_a R_a e^{\epsilon'R_a}\sup_I p_a,
\]
which is finite because
$R e^{\epsilon'R}\le C_{\epsilon,\epsilon'}e^{\epsilon R}$.
Artificial insertions cancel before taking total variation.  The removed
high-homogeneity pieces satisfy the common two-sided recovery moment and are
summed by Proposition~\ref{prop:recovery-window-summability}.
\end{proof}

The following consequence records the exact countable-IFS interface rather
than hiding it in a finite-horseshoe approximation.

\begin{proposition}[Finite-difference countable-IFS bridge]
\label{prop:countable-ifs-bridge}
Let the central unstable quotient be a one-state full-branch
$C^{1+\alpha}$ IFS with uniform inverse contraction, bounded distortion, a
normalized Gibbs potential of summable variations, and
\eqref{eq:lps-light-tail}.  For words $\mathbf a,\mathbf b$ of length $N$
and $x,y\in I$, put
\[
 \Delta_N((\mathbf a,x),(\mathbf b,y))
 =\{\log|h_{\mathbf a}'(x)|-\log|h_{\mathbf b}'(x)|\}
  -\{\log|h_{\mathbf a}'(y)|-\log|h_{\mathbf b}'(y)|\}.
\]
Assume the LPS quantitative nonlinearity estimate: for some
$C>0$, $\Theta>0$, and $\rho\in(0,1)$,
\begin{equation}
\label{eq:finite-difference-qnl}
 \sum_{\mathbf a,\mathbf b,\mathbf c,\mathbf d\in\mathcal A^N}
 p_{\mathbf a}p_{\mathbf b}p_{\mathbf c}p_{\mathbf d}
 \one_{J}\!\left(
 \Delta_N((\mathbf a,x_{\mathbf c}),
          (\mathbf b,x_{\mathbf d}))\right)
 \le C\{|J|^\Theta+\rho^N\}
\end{equation}
for every interval $J\subset\mathbb R$; Gibbs weights may be evaluated at
arbitrary reference points in their cylinders.  Then the induced SRB Gibbs
measure has the weighted nonstationary-phase Fourier estimate of
Leclerc--Paukkonen--Sahlsten for every phase with the stated
$C^\alpha$ log-derivative bounds.  The conclusion remains valid for the
marked physical face amplitudes of
Lemma~\ref{lem:weighted-face-qnl-transfer}.
\end{proposition}

\begin{proof}
Equation~\eqref{eq:finite-difference-qnl} is precisely the QNL hypothesis in
Theorem~1.1 of \cite{LPS}.  That theorem is formulated for finite or
countable $C^{1+\alpha}$ systems and H\"older amplitudes; it neither
differentiates $\log h_a'$ nor assumes a $C^2$ holonomy.  The light-tail
moment controls the countable alphabet.  The weighted-face conclusion follows
from Lemma~\ref{lem:weighted-face-qnl-transfer} and the summable
scale-normalized amplitude norm \eqref{eq:scaled-amplitude-norm}.
\end{proof}

The former $C^2$ multiplier/UNI route is deliberately not used.  For a
billiard quotient the stable holonomy is generally only $C^{1+\alpha}$, so
the derivative of $\log h_a'$ need not exist.  The correct gate input is the
finite-difference estimate \eqref{eq:finite-difference-qnl}.

\begin{proposition}[Clean homoclinic horseshoe supplies a finite QNL gate]
\label{prop:homoclinic-horseshoe-uni}
Let $p$ be the regular fixed point of the squared billiard map associated with
the normal period-two orbit in Proposition
\ref{prop:exact-circular-anosov}.  If its Anosov cocycle satisfies
$\mathfrak A(p)\ne0$, then there is a compact finite clean horseshoe
$\Omega_p$, contained in regular branches and containing $p$, such that in at
least one time orientation its temporal distance satisfies QNL for every
H\"older equilibrium state.  Thus $\Omega_p$ supplies a finite topological
      QNL seed; Proposition~\ref{prop:context-polynomial-blowup} transfers its
      polynomial escape to the physical collision-SRB bracket.
\end{proposition}

\begin{proof}
Lemma~5.1 of Lima--Obata--Poletti~\cite{LOP} says that finite-horizon
dispersing billiards have a single homoclinic class at regular hyperbolic
points.  Choose a second regular hyperbolic periodic point $q$.  The theorem
gives clean transverse intersections
$W^u(p)\pitchfork W^s(q)$ and $W^u(q)\pitchfork W^s(p)$.  The inclination
lemma then gives a nontrivial clean transverse homoclinic point of $p$; this
avoids the vacuous substitution $x=y=p$.  The Smale--Birkhoff construction, using
a finite segment of this clean homoclinic orbit, gives a compact finite
horseshoe $\Omega_p$ on which the billiard map is smooth.  In symplectic
Birkhoff coordinates the map is area preserving; the local dynamics near
$\Omega_p$ may be viewed as a smooth Axiom--A basic set.

By the contrapositive of the Anosov cocycle obstruction, together with
Theorem~4.3 of Leclerc~\cite{Leclerc2026}, in one of the two time orientations
the temporal distance on $\Omega_p$ satisfies QNL for every H\"older
equilibrium state, in particular for the equilibrium state of the geometric
potential restricted to this finite subsystem.  No comparison of that
horseshoe measure with collision SRB is made here.
\end{proof}

There is a regularity caveat.  The stable quotient of a smooth horseshoe is in
general only $C^{1+\alpha}$, precisely because the stable holonomy need not be
$C^2$.  Therefore the finite QNL conclusion above cannot be fed directly
into the $C^2$ non-linearizability criterion of
Algom--Rodriguez Hertz--Wang.  The route below instead propagates the
physical gate's finite-difference gap to the countable
$C^{1+\alpha}$ inducing scheme.  The LPS theorem is recorded as a background
alternative; it is not invoked in the principal fractional-phase proof.

The light-tail hypothesis has a concrete billiard reduction.  If a branch
$a$ has return time $R_a$ and grazing rank $k_a$, bounded distortion gives,
schematically,
\[
 |h_a'|^{-\gamma_0}\le
 C e^{c\gamma_0R_a}(1+k_a)^{q\gamma_0}.
\]
Thus an exponential return moment and any strictly positive polynomial SRB
tail in $k_a$ imply \eqref{eq:lps-light-tail} after choosing
$\gamma_0>0$ sufficiently small.  This is strictly weaker than uniform
control of every rank-$L$ smooth extension.

\begin{lemma}[Return and grazing moments imply the LPS light tail]
\label{lem:billiard-light-tail}
Assume the normalized branch weights $p_a(x)=e^{\phi(h_ax)}$ satisfy
\[
 \sum_a\sup_x p_a(x)e^{\epsilon R_a}(1+k_a)^\delta<\infty
\]
for some $\epsilon,\delta>0$, and suppose
$|h_a'|^{-1}\le Ce^{cR_a}(1+k_a)^q$.
Then \eqref{eq:lps-light-tail} holds for every
\[
 0<\gamma_0<\min\{\epsilon/c,\delta/q\}.
\]
\end{lemma}

\begin{proof}
Raise the inverse-derivative estimate to $\gamma_0$ and absorb the two
factors into the assumed exponential and polynomial moment.
\end{proof}

There is also a pressure-based way to obtain the required moment.  The
geometric-potential theory for finite-horizon dispersing billiards has a
spectral gap for $-t\log J^uT$ on every compact subinterval of
$(0,t_*)$, where $t_*>1$, and its modified one-step expansion is uniform for
$t\ge t_0>0$~\cite{DLreview}.  Since the SRB value is $t=1$, the LPS factor
$|h_a'|^{-\gamma_0}$ merely changes the geometric exponent from $1$ to
$1-\gamma_0$.

\begin{proposition}[Geometric-pressure route to the light tail]
\label{prop:pressure-light-tail}
Let a homogeneous full-branch inducing scheme have an exponential return
tail at the SRB parameter.  Assume its induced distortion is uniform and
that the geometric transfer estimates hold for
$t\in[1-\varepsilon,1]$.  Then, after decreasing
$\varepsilon>0$ if necessary, its normalized induced SRB potential satisfies
\eqref{eq:lps-light-tail} for every sufficiently small
$\gamma_0>0$.
\end{proposition}

\begin{proof}
Group branches by return time $R_a=n$.  Up to the uniform distortion
coboundary, the summand in \eqref{eq:lps-light-tail} is the inverse unstable
Jacobian to the power $1-\gamma_0$.  The geometric transfer estimate at
$t=1-p\gamma_0$ bounds the corresponding $p$-moment by
$Ce^{c(p\gamma_0)n}$, with $c(p\gamma_0)\to0$ as
$\gamma_0\to0$.  H\"older's inequality combines this with the SRB return
tail $Ce^{-an}$ to give
\[
 \sum_{R_a=n}\sup_x
 e^{\phi(h_ax)}|h_a'(x)|^{-\gamma_0}
 \le C e^{-a n/q+c(p\gamma_0)n/p}.
\]
Choose $p,q>1$ conjugate and then $\gamma_0$ so that the exponent is
negative; summing over $n$ proves the claim.  Homogeneity-strip subdivision
is already included in the geometric one-step expansion.
\end{proof}

\begin{lemma}[Fixed-gate hitting and two-context moments]
\label{lem:fixed-gate-context-moment}
Let $F$ be a one-state full-branch Gibbs--Markov map and let $\mathcal G$ be a
fixed finite gate cylinder of return depth $L$.  Suppose bounded distortion
gives
\begin{equation}
\label{eq:fixed-gate-probability}
 \inf_C\mu^u(\mathcal G\mid C)\ge\eta_G>0
\end{equation}
for every past cylinder $C$, and suppose
\begin{equation}
\label{eq:return-exponential-moment}
 \mathbb E e^{\epsilon R}<\infty
\end{equation}
for some $\epsilon>0$.  If $N_G$
is the first gate-hitting time in $L$-blocks and $C^+$ is the total collision
length of the future entrance context, then there is $\epsilon_G>0$ such that
\[
 \mu^u(N_G>k)\le(1-\eta_G)^k,
 \qquad
 \mathbb E e^{\epsilon_G C^+}<\infty.
\]
The reverse quotient gives the same estimate for $C^-$.  On the natural
extension no independence is needed:
\begin{equation}
\label{eq:two-context-moment}
 \mathbb E e^{\epsilon_G'(C^-+C^+)}<\infty
\end{equation}
for every sufficiently small $\epsilon_G'>0$.
\end{lemma}

\begin{proof}
Condition \eqref{eq:fixed-gate-probability} and the Markov property give
$\mu^u(N_G>k+1\mid N_G>k)\le1-\eta_G$, proving the geometric tail.  Let
$K(\theta)$ be the uniform conditional exponential moment of one return;
$K(0)=1$ and $K(\theta)<\infty$ for small $\theta>0$.  Bounded Gibbs
distortion gives
\[
 \mathbb E\!\left[e^{\theta C^+};N_G=k\right]
 \le C\{(1-\eta_G)K(\theta)^L\}^{k-1}K(\theta)^L.
\]
Choose $\theta$ so that the quantity in braces is smaller than one and sum in
$k$.  Time reversal proves the past estimate.  Finally
\[
 \mathbb P(C^-+C^+>t)
 \le\mathbb P(C^->t/2)+\mathbb P(C^+>t/2),
\]
so the two marginal exponential tails imply
\eqref{eq:two-context-moment} without independence.
\end{proof}

For the marked quotient, choose one fixed finite gate word.  The Gibbs ratio
of that word after an arbitrary prefix differs from its unconditional weight
by at most the uniform distortion factor.  Hence
\eqref{eq:fixed-gate-probability} holds with
$\eta_G=C_D^{-1}\mu^u(\mathcal G)>0$; it is not an additional mixing
assumption.

\begin{lemma}[Joint return--maximal-grazing moment]
\label{lem:homogeneity-grazing-moment}
Let $Y$ be the base of an exponential Young return, $R$ its return time, and
let $B(x)$ be the largest homogeneity rank visited before $R(x)$.  In
$(r,p)=(r,\sin\varphi)$ coordinates suppose
\[
 \mu\{1-|p|\lesssim2^{-2b}\}\le C2^{-2b},
 \qquad \mu_Y(R>L)\le Ce^{-cL}.
\]
Then $B$ has an exponential tail.  More precisely, for every fixed $A>0$,
\begin{equation}
\label{eq:maximal-grazing-tail}
 \mu_Y(B\ge b)
 \le Ce^{-cAb}+CAb\,2^{-2b}.
\end{equation}
Consequently, for sufficiently small $\epsilon,\kappa>0$,
\begin{equation}
\label{eq:small-grazing-moment}
 \mathbb E_Y e^{\epsilon R+\kappa B}<\infty.
\end{equation}
If the inverse-Jacobian and coarea loss of a return word is bounded by
$Ce^{c_0R+qB}$, then \eqref{eq:lps-light-tail} and the marked-amplitude
moment hold after decreasing the exponents.  No independence of $R$ and $B$
is required.
\end{lemma}

\begin{proof}
On $\{R\le Ab\}$ take a union bound over the at most $Ab$ collision times.
Since $\mu_Y(E)\le\mu(E)/\mu(Y)$ for every event $E$ and the collision SRB
measure is invariant, each visit to a rank-$b$ boundary layer costs at most
$C2^{-2b}$.  Adding $\mu_Y(R>Ab)$ proves
\eqref{eq:maximal-grazing-tail}.  Choose $A$ so that both terms decay
exponentially in $b$; hence $\mathbb E e^{q\kappa B}<\infty$ for some
$q>1$ and small $\kappa$.  For conjugate $p,q>1$, H\"older gives
\[
 \mathbb E e^{\epsilon R+\kappa B}
 \le (\mathbb E e^{p\epsilon R})^{1/p}
      (\mathbb E e^{q\kappa B})^{1/q}<\infty.
\]
Absorb $c_0R+qB$ with smaller exponents and apply
Lemma~\ref{lem:billiard-light-tail}.  The same argument applies to the
scale-normalized coarea amplitude.
\end{proof}

Lemma~\ref{lem:homogeneity-grazing-moment} verifies the grazing factor from
the collision density and strip width, while
Proposition~\ref{prop:central-marked-moment} supplies the moving-face
amplitude budget on the same quotient.  Lemma
\ref{lem:fixed-gate-context-moment} supplies the entrance/exit moment.  Thus
the particular full-branch quotient uses one joint return/grazing/context
ledger rather than an unverified appeal to the literature theorem.

This route also exposes a flaw in the finite-horseshoe detour.  A conservative
pasting lemma~\cite{Teixeira} can locally modify or smooth a volume-preserving map, but it
does not automatically extend a finite family of billiard branches to a
global area-preserving Axiom--A diffeomorphism whose equilibrium state has the
required collision-SRB weights.  Even a successful extension can put a
fractal Gibbs conditional on a saddle basic set in place of the absolutely
continuous unstable conditional of the billiard SRB measure.  The direct
countable-IFS quotient preserves the correct measure and should therefore be
the primary route.

Thus light-tail is not expected to be a new principal theorem once the
standard homogeneous inducing scheme is put in the one-state form required
above.  The exact value $\mathfrak A=-1$, the clean homoclinic theorem, and
Proposition~\ref{prop:homoclinic-horseshoe-uni} produce the clean polynomial
escape seed without naming a connector.  Proposition
\ref{prop:physical-legal-child} records the local physical slope transfer, and
the cohomological bracket below identifies its physical collision-SRB object.

\subsection{The cohomological bracket and polynomial context blow-ups}

The quotient regularity is only $C^{1+\alpha}$.  The useful exact object is
therefore a four-point difference of Birkhoff sums, not a derivative of a
holonomy.  Let $f$ denote the billiard map restricted to the finite clean
basic set, let $F$ be its stable factor, and put
$\tau_f=\log J^uf$ and $\tau_F=\log|F'|$.  As in
\cite[Lemma~3.15]{Leclerc2026}, choose a H\"older function $\theta$ such that
\begin{equation}
\label{eq:potential-cohomology}
 \tau_f=\tau_F\circ\pi+\theta\circ f-\theta .
\end{equation}
For $p,q$ in one Markov rectangle define
\begin{equation}
\label{eq:physical-temporal-bracket}
 \Delta(p,q)=\sum_{j\in\mathbb Z}
 \{\tau_f(f^jp)-\tau_f(f^j[p,q])
   -\tau_f(f^j[q,p])+\tau_f(f^jq)\}.
\end{equation}

\begin{lemma}[Cohomological bracket identification]
\label{lem:cohomological-bracket-identification}
Let $\mathbf a,\mathbf b$ be the length-$N$ past words of $p,q$ and let
$x=\pi p$, $y=\pi q$.  With inverse branches $g_{\mathbf a},g_{\mathbf b}$
of $F$, set
\begin{align}
 \mathcal D_N(\mathbf a,\mathbf b;x,y)
  :={}&S_N\tau_F(g_{\mathbf a}x)-S_N\tau_F(g_{\mathbf a}y)\notag\\
      &-S_N\tau_F(g_{\mathbf b}x)+S_N\tau_F(g_{\mathbf b}y).
\label{eq:factor-four-bracket}
\end{align}
Then
\begin{equation}
\label{eq:bracket-tail}
 |\mathcal D_N(\mathbf a,\mathbf b;x,y)-\Delta(p,q)|
 \le C\kappa^{\alpha N}.
\end{equation}
 The constants are uniform on the persistent finite basic set and for
 $|s|<s_0$.  Entrance and exit contexts are not included in this statement;
 their finite jets and holonomy tails are treated explicitly in Proposition
 \ref{prop:context-polynomial-blowup}.
\end{lemma}

\begin{proof}
Insert \eqref{eq:potential-cohomology} in
\eqref{eq:physical-temporal-bracket}.  The four $\theta$-terms telescope.
For $j\ge0$, stable projection gives
$\pi(f^jp)=\pi(f^j[p,q])$ and
$\pi(f^jq)=\pi(f^j[q,p])$, so the positive-time factor terms cancel
exactly.  Writing the negative iterates through the past inverse words gives
\eqref{eq:factor-four-bracket}.  The omitted tail starts at $N+1$; H\"older
regularity of $\tau_F$ and inverse contraction give
\[
 \sum_{j>N}C\kappa^{\alpha j}\le C'\kappa^{\alpha N}.
\]
 This is precisely the proof of \cite[Lemma~3.15]{Leclerc2026} on the clean
 basic set.  Hyperbolic continuation makes the constants uniform in $s$.
\end{proof}

The next statement transfers the all-scale escape in
\cite[Propositions~4.31--4.32]{Leclerc2026} from the horseshoe conditional to
the physical collision-SRB conditional.  This is the only support transfer
that is required.

The preceding bracket must still be matched to the physical endpoint maps.
For a stable projection $\pi_s$ on a clean rectangle, the stable holonomy
Jacobian satisfies the convergent identity
\begin{equation}
\label{eq:exact-holonomy-series}
 \log J\pi_s(r)=\sum_{j\ge0}
 \{\tau_f(f^jr)-\tau_f(f^j\pi_s r)\}.
\end{equation}
This is the Jacobian formula of \cite[Lemma~C.1]{Leclerc2026}.  In particular,
entrance and exit holonomies are not discarded as unspecified coboundaries:
their four endpoint series are inserted into \eqref{eq:physical-temporal-bracket};
the common terms cancel and the unpaired tail is bounded by
$C\kappa^{\alpha N}$.

\begin{proposition}[Context--polynomial blow-up]
\label{prop:context-polynomial-blowup}
Assume $\mathfrak A(p)=-1$ and let $I_c=t_c+\ell_c[-1,1]$ be any homogeneous
resolving cylinder whose central excursion enters the persistent clean basic
set.  For the full reduced-slope bracket $S_c$, including the two endpoint
holonomy Jacobians, let $N$ be the total central resolving depth and let
$L_c=|c^-|+|c^+|$ be the entrance/exit length.  There are a fixed degree $d$,
a rough template $\mathcal X_{\omega_c}$ from the clean basic set, a
polynomial $P_c$ of degree at most $d$, and a context cost $H_c\ge1$ such that
\begin{equation}
\label{eq:context-polynomial-blowup}
 \ell_c^{-1}\{S_c(t_c+\ell_c u)-S_c(t_c)\}
 =\mathcal X_{\omega_c}(u)+P_c(u)+E_c(u),\qquad
 \|E_c\|_\infty\le C\left(H_c\ell_c^\alpha+
       \frac{\kappa^{\alpha N}}{\ell_c}\right).
\end{equation}
Moreover $H_c\le Ce^{aL_c}$, and for some $p>1$ the context ledger gives
\begin{equation}
\label{eq:context-jet-moment}
 \sup_{|s|<s_0}\sum_c w_c H_c^p<\infty .
\end{equation}
The degree and the escape constants of $\mathcal X_{\omega_c}$ modulo
$\mathbb R_d[X]$ are independent of $c$ and $s$.  In particular, if
\begin{equation}
\label{eq:resolving-depth-scale}
 N\ge q_*|\log\ell_c|,
 \qquad q_*>\frac{1+\alpha}{\alpha|\log\kappa|},
\end{equation}
then the second error in \eqref{eq:context-polynomial-blowup} is
$O(\ell_c^\alpha)$.
\end{proposition}

\begin{proof}
Use \eqref{eq:exact-holonomy-series} for both endpoint maps and split each
series at the entrance and exit times.  The four common tails cancel in the
bracket.  Each remaining finite clean context sum is $C^{1+\alpha}$ on its
homogeneous component.  Its affine Taylor polynomial is linear (and is
regarded as an element of $\mathbb R_d[X]$); the normalized remainder is
bounded by its $C^{1+\alpha}$ norm times $\ell_c^\alpha$.  The chain rule and
uniform one-step distortion give $H_c\le Ce^{aL_c}$.  Choose $p>1$ so that
$pa$ is smaller than the exponent in \eqref{eq:two-context-moment}; this proves
\eqref{eq:context-jet-moment}.  Lemma
\ref{lem:cohomological-bracket-identification} identifies the central part
with the factor four-point cocycle with error $O(\kappa^{\alpha N})$.
Condition \eqref{eq:resolving-depth-scale}, rather than an unrecorded fixed
acceleration, makes its normalized value $O(\ell_c^\alpha)$.  Finally
\cite[Propositions~4.31--4.32]{Leclerc2026}, activated by
$\mathfrak A(p)\ne0$, gives uniform all-location, all-scale escape of the rough
template from every polynomial of degree at most $d$.  No fixed acceleration
is used.
\end{proof}

\subsection{Native-coordinate single-source slope energy}

The following two conditions isolate the genuinely new clean-gate geometry
on the \emph{actual} collision--SRB interval.  They are not inferred from a
smooth extension of a fractal horseshoe conditional.

\begin{definition}[Single-source bracket condition $\mathrm{SS}$]
\label{def:single-source-condition}
Fix one orientation and condition on the opposite symbolic record, the value
$t'$, and the two finite contexts.  For a native interval $Q$ write
$z=\Phi_{\beta,Q}(t)$ and
$\sigma_Q=\operatorname{diam}\Phi_{\beta,Q}(Q)$.  Put
$d_Z=\max\{1,d\}$, where $d$ is
the fixed degree in Proposition~\ref{prop:context-polynomial-blowup}.
Condition $\mathrm{SS}(d_Z)$ means that
\begin{equation}
\label{eq:exact-reduced-slope}
 S_\beta(t)-S_\beta(t')
 =\ell_{\beta,Q}\{X^{[N]}_{\beta,Q}(z)-P_{\beta,Q,t'}(z)\}
  +\mathcal R_{\beta,Q,t'}(z),
\end{equation}
where $P_{\beta,Q,t'}\in\mathbb R_{d_Z}[z]$ and
$X^{[N]}_{\beta,Q}$ is the finite endpoint series on the actual depth-$N$
collision cylinder.  On the Cantor core it is within $C\kappa^N$ of one
native Leclerc rough template, and
\begin{equation}
\label{eq:single-source-coefficient}
 0<c_\ell\le|\ell_{\beta,Q}|\le C_\ell,
 \qquad
 \|\mathcal R_{\beta,Q,t'}\|_\infty
 \le C\{H_\beta\sigma_Q^{1+\alpha_B}+\kappa^N\}.
\end{equation}
The endpoint identity proving that every other $t$-dependent rough term is
absent or belongs to the fixed-degree polynomial is part of the condition;
freezing a symbolic name alone is not enough.
\end{definition}

\begin{definition}[Finite-depth physical thickening $\mathrm{FT}$]
\label{def:physical-legal-extension}
Let $\mu^u_{\beta,Q}$ be the normalized collision--SRB conditional on a
native interval $Q$.  The persistent clean gate satisfies
$\mathrm{FT}(d_Z)$ when the following finite-time statements hold.
\begin{enumerate}[label=\textup{(FT\arabic*)}]
\item The clean return branches defining the Axiom--A alphabet extend to
      open $C^{1+\alpha_B}$ full branches on physical unstable intervals.
      Every finite legal word $a_1\cdots a_J$ has an open collision cylinder
      $V_{a_1\cdots a_J}$, and the stable holonomy from the physical gate to
      the reference unstable interval has Jacobian bounded above and below
      with uniform scaled log--H\"older distortion.
\item Leclerc's equilibrium state is chosen for the induced geometric
      potential $-\log J^uF$ (with the normalizing coboundary).  Consequently
      the equilibrium weight of every depth-$J$ native cylinder and the
      relative collision--SRB mass of its open thickening are uniformly
      comparable.  The comparison is required only for finite words and is
      uniform in $J$ by the Gibbs distortion bound.
      Uniformly for every $P\in\mathbb R_{d_Z}[z]$, the finite cutting
      antichain splits into bad cylinders of total Gibbs weight at most
      $C\sigma_Q^\gamma$ and good inner cylinders satisfying the robust core
      margin $|X_{\beta,Q}-P|\ge4\sigma_Q^{1+\alpha_L}$.
\item On $V_{a_1\cdots a_J}$ the finite endpoint/holonomy series
      $X^{[J]}_{\beta,Q}$ is defined on the full physical cylinder and
      \begin{equation}
      \label{eq:finite-thickening-shadow-error}
       \|X^{[J]}_{\beta,Q}-X_{\beta,Q}\|_\infty\le C\kappa^J
      \end{equation}
      on the Cantor core.  There is a fixed $d_{\rm ph}\ge0$ such that the
      context polynomials occurring in $\mathrm{SS}(d_Z)$ are admissible in
      the sense that, on a good parent,
      \begin{equation}
      \label{eq:physical-lipschitz-mark}
       \Lambda_V(P):=1+\operatorname{Lip}_V(X^{[J]}_{\beta,Q}-P)
       \le CH_\beta\sigma_Q^{-d_{\rm ph}}.
      \end{equation}
      Every good symbolic child is continued by pure Markov refinement until
      each retained open physical descendant $W\subset V$ has
      \begin{equation}
      \label{eq:two-scale-physical-diameter}
       \operatorname{diam}W\le
       \rho_V:=\frac{\sigma_Q^{1+\alpha_L}}{8\Lambda_V(P)}.
      \end{equation}
      Every such descendant meets the corresponding Cantor child, all legal
      descendants are retained, their physical and Gibbs weights remain
      uniformly comparable, and the added depth is at most
      $C(1+d_{\rm ph}+d_0)|\log\sigma_Q|$ on the good contexts of
      Lemma~\ref{lem:gate-buffer}.
      This is the second physical scale; no affine jet is absorbed into $P$.
\item The finite-time coupling/growth construction returns a subfamily of
      actual conditional mass at least $\eta_{\rm reg}>0$ to the clean gate.
      Failed trials obey
      $L_{j+1}\le(1-\eta_{\rm reg})L_j+Ce^{-c(j+|\log\sigma_Q|)}$;
      endpoint overshoots have bounded multiplicity, all continuations are
      retained, and $J$ trials plus a depth-$J$ word cost
      $C(J+|\log\sigma_Q|)$ physical iterates.
      More precisely, after the second-scale refinement, the stopping
      antichain has disjoint open thickenings $W_\omega$ and a measurable
      physical gap $G_J$ such that
      \begin{equation}
      \label{eq:finite-antichain-physical-cover}
       Q=G_J\,\dot\cup\!\mathop{\dot\bigcup}_{\omega\in\mathcal A_J}W_\omega
       \quad(\bmod\ \mu^u_{\beta,Q}),\qquad
       \mu^u_{\beta,Q}(G_J)\le Ce^{-cJ}.
      \end{equation}
      Endpoint overshoots are assigned to their first failed ancestor, so
      none is omitted or counted twice.
\end{enumerate}
No assertion in $\mathrm{FT}$ says that a polynomial sublevel already loses a
fixed amount of actual SRB mass.  However, \textup{(FT3)--(FT4)} are primitive
actual-SRB physical-thickening inputs: in particular, Leclerc's cutting theorem
does not supply the second-scale open cover, the Gibbs/physical comparison, or
the exponentially small physical gap.  The conclusion below is the logical
combination of those inputs with the finite Leclerc margin.  No infinite
Cantor conditional is pushed through a smooth map and declared absolutely
continuous.
\end{definition}

\begin{proposition}[Single-source physical bracket]
\label{prop:exact-reduced-slope}
Under $\mathrm{SS}(d_Z)$,
\eqref{eq:exact-reduced-slope}--\eqref{eq:single-source-coefficient} hold on
every labelled native interval.  In the dominant coordinates there are
positive labelled scales $a_X,a_Y$ such that
\begin{equation}
\label{eq:dominant-coordinate-normalization}
 c\le |X'|/a_X,|Y'|/a_Y\le C,
 \qquad \Xi=|\xi_1|a_X\asymp|\xi_2|a_Y.
\end{equation}
\end{proposition}

\begin{proof}
The bracket assertion is precisely $\mathrm{SS}(d_Z)$.  The last display is
the nonvanishing part of the clean-gate chart condition after contracted
coordinates are removed.  No uniform lower bound is asserted: the inverse
scale cost is the explicit factor $\mathfrak d_\lambda$ in
\eqref{eq:event-physical-norm}.
\end{proof}

\begin{lemma}[Gate, buffer, and native scale]
\label{lem:gate-buffer}
Let $\alpha_L<\alpha_B$, put
\begin{equation}
\label{eq:native-energy-scale}
 \sigma=(r/c_\ell)^{1/(1+\alpha_L)},\qquad
 d_0=\frac{\alpha_B-\alpha_L}{2},
\end{equation}
and call a context good when $H_\beta\le\sigma^{-d_0}$.  Choose
$N=\lceil C_{\rm buf}|\log r|\rceil$ so that
$\kappa^N\le r/10$.  Then on good contexts the error in
\eqref{eq:single-source-coefficient} is $o(r)$, while
\begin{equation}
\label{eq:bad-context-native-scale}
 \mathbb P\{H_\beta>\sigma^{-d_0}\}
 \le \mathbb E(H_\beta^p)\,r^{pd_0/(1+\alpha_L)}.
\end{equation}
The regeneration/cutting time from $\mathrm{FT}$ plus $N$ is
$O(|\log r|)$ and all buffer continuations are retained, so the refinement
preserves conditional mass.
\end{lemma}

\begin{proof}
On a good context,
$H_\beta\sigma^{1+\alpha_B}
\le\sigma^{1+(\alpha_B+\alpha_L)/2}=o(\sigma^{1+\alpha_L})=o(r)$.
Markov's inequality gives \eqref{eq:bad-context-native-scale}.  The remaining
claims follow from exponential gate hitting and the mass-preserving buffer.
\end{proof}

\begin{proposition}[Finite-depth thickening implies physical polynomial escape]
\label{prop:physical-legal-child}
Under $\mathrm{FT}(d_Z)$ there is $\gamma_{\rm ph}>0$ such that,
uniformly for the admissible context polynomials
$P\in\mathbb R_{d_Z}[z]$ satisfying
\eqref{eq:physical-lipschitz-mark}, every native interval obeys
\begin{equation}
\label{eq:local-clean-sublevel}
 \mu^u_{\beta,Q}
 \{|X^{[J]}_{\beta,Q}-P|\le\sigma_Q^{1+\alpha_L}\}
 \le C\sigma_Q^{\gamma_{\rm ph}}.
\end{equation}
The same bound holds after multiplication by a positive marked amplitude,
with the corresponding scaled amplitude mark.
\end{proposition}

\begin{proof}
Put $J=\lfloor c_0|\log\sigma_Q|\rfloor$ with $c_0$ large enough that
$\kappa^J\le\sigma_Q^{1+\alpha_L}/100$.  Apply Lemma~4.34 of
\cite{Leclerc2026} only to its finite depth-$J$ stopping antichain.  Apply
Proposition~4.33 at the enlarged threshold
$4\sigma_Q^{1+\alpha_L}$.  The union of antichain cylinders which can meet
that sublevel has equilibrium weight at most $C\sigma_Q^{\gamma}$, while
every remaining inner cylinder has the robust margin
\begin{equation}
\label{eq:robust-finite-child-margin}
 |X_{\beta,Q}-P|\ge2\sigma_Q^{1+\alpha_L}
\end{equation}
on its Cantor core.  This is a finite statement: no infinite intersection is
taken.

Use \textup{(FT1)} to thicken the whole stopping antichain to disjoint open
physical collision cylinders.  Since the equilibrium and actual conditional
weights are generated by the same normalized geometric potential,
\textup{(FT2)} compares the sum of the physical weights of the potentially
bad cylinders with $C\sigma_Q^\gamma$.  The shadow error
\eqref{eq:finite-thickening-shadow-error} and the choice of $J$ first give
$|X^{[J]}_{\beta,Q}-P|\ge(3-o(1))\sigma_Q^{1+\alpha_L}$ at every Cantor point
of a good child.  Now use the second refinement in \textup{(FT3)}.  If
$z_W$ is a Cantor point in a retained descendant $W$, then for every $z\in W$,
\[
 |(X^{[J]}-P)(z)-(X^{[J]}-P)(z_W)|
 \le\Lambda_V(P)\operatorname{diam}W
 \le\tfrac18\sigma_Q^{1+\alpha_L}.
\]
Thus the original event, with the original polynomial $P$, has
$|X^{[J]}_{\beta,Q}-P|>\sigma_Q^{1+\alpha_L}$ on the whole open descendant.
No change of polynomial is made.  The physical cover
\eqref{eq:finite-antichain-physical-cover} and the first-ancestor overshoot
assignment give total uncovered mass $Ce^{-cJ}$ without a hidden covering
step.  Hence
\[
 \mu^u_{\beta,Q}
 \{|X^{[J]}_{\beta,Q}-P|\le\sigma_Q^{1+\alpha_L}\}
 \le C\sigma_Q^\gamma+Ce^{-cJ}
 \le C\sigma_Q^{\gamma_{\rm ph}}
\]
for a smaller $\gamma_{\rm ph}>0$.  Multiplication by a positive amplitude is
made before normalization and costs its scaled supremum.  This proves
\eqref{eq:local-clean-sublevel} on the actual collision conditional using
only finite cylinders.
\end{proof}

\begin{theorem}[Physical single-source energy tree]
\label{thm:direct-slope-small-ball}
Let $\mu_\beta$ be the domain conditional and
$\nu_\beta=(S_\beta)_*\mu_\beta$.  With
\begin{equation}
\label{eq:physical-energy-exponent}
 \Theta_E=\frac{\min\{\gamma_{\rm ph},pd_0\}}{1+\alpha_L}>0,
\end{equation}
one has
\begin{align}
 \sum_\beta w_\beta\iint
   \mathbf 1_{\{|S_\beta(t)-S_\beta(t')|\le r\}}
   \,d\mu_\beta(t)d\mu_\beta(t')
 &\le Cr^{\Theta_E}+Ce^{-cN},\label{eq:conditional-energy-bound}\\
 \sup_z\sum_\beta w_\beta\nu_\beta([z-r,z+r])
 &\le Cr^{\Theta_E/2}+Ce^{-cN}.
 \label{eq:slope-value-small-ball}
\end{align}
\end{theorem}

\begin{proof}
Fix $t'$ and all frozen physical data.  By
\eqref{eq:single-source-coefficient}, the event
$|S_\beta(t)-S_\beta(t')|\le r$ is contained, after enlarging constants, in
 $|X^{[N]}_{\beta,Q}-P_{\beta,Q,t'}|\le\sigma^{1+\alpha_L}$ with
$\sigma=(r/c_\ell)^{1/(1+\alpha_L)}$.  Enlarge $C_{\rm buf}$ once so that
$N$ is at least the stopping depth $J$ in Proposition
\ref{prop:physical-legal-child}.  The additional endpoint-series tail from
$J$ to $N$ is $O(\kappa^J)$ and is absorbed by the enlarged sublevel
threshold.  On good contexts combine
Proposition~\ref{prop:exact-reduced-slope}, Lemma~\ref{lem:gate-buffer}, and
\eqref{eq:local-clean-sublevel}.  Equation \eqref{eq:bad-context-native-scale}
handles bad contexts, and the buffer/gate tails give $e^{-cN}$.  Integrate in
$t'$ and then over branches; this order is essential because the bad-context
estimate is averaged over the physical branch law and is not uniform for a
fixed $\beta$.  Finally use
\[
 \nu_\beta(I)^2\le
 (\nu_\beta\otimes\nu_\beta)\{|u-v|\le2|I|\}
\]
and Jensen in the positive weights $w_\beta$ to obtain
\eqref{eq:slope-value-small-ball}.
\end{proof}

\begin{corollary}[Central slope small-ball]
\label{cor:central-g2}
The complete central collision--SRB quotient satisfies
\eqref{eq:slope-value-small-ball} with constants uniform in $s$.
\end{corollary}

\begin{proof}
Apply Theorem~\ref{thm:direct-slope-small-ball} and the uniform marked
return/grazing/context moments.
\end{proof}

\begin{lemma}[Slope small-ball controls the balanced phase]
\label{lem:slope-good-bad-phase}
For $\psi_\xi=\xi_1X+\xi_2Y$ on an affine-normalized gate chart, let
\begin{equation}
\label{eq:balanced-physical-frequency}
 \Xi=|\xi_1|a_X\asymp|\xi_2|a_Y,
\end{equation}
with $a_X,a_Y$ as in \eqref{eq:dominant-coordinate-normalization}.  The mass of
$\{|\psi_\xi'|\le\Xi r\}$ is at most
 $C(r^{\Theta/2}+e^{-cn})$.  Smooth value-space cutoffs on its complement have
$C^\alpha$ norm $O(r^{-1})$ and satisfy
$\inf|\psi_\xi'|\ge c\Xi r$.
\end{lemma}

\begin{proof}
By Proposition~\ref{prop:exact-reduced-slope},
\[
 \psi_\xi'=X'(\xi_1+\varepsilon_{\rm or}\xi_2e^S).
\]
The lower and upper bounds in
\eqref{eq:dominant-coordinate-normalization} convert the bracketed factor to
the physical scale \eqref{eq:balanced-physical-frequency}.  Only
opposite-oriented coefficients can resonate, and the resonant set is
contained in one interval $|S-z_\xi|\le Cr$.  Apply
\eqref{eq:slope-value-small-ball}.  Composing a fixed smooth three-term
partition with
$(\xi_1+\varepsilon_{\rm or}\xi_2e^S)/(\Xi r)$ gives the cutoff and norm
bounds.
\end{proof}
\begin{lemma}[Typed global double projection before disintegration]
\label{lem:global-label-recentering}
Let $J$ be the complete finite signed response measure on $M\times M$ in
common physical coordinates, let $m_J^\pm$ be its marginals, and let
$\nu^\pm$ be the two fixed smooth probability references.  Define the
current-side projection
\begin{equation}
\label{eq:labelwise-recentering}
 \mathsf D_{\rm cur}J=:J^\circ
 =J-m_J^-\otimes\nu^+-\nu^-\otimes m_J^+
   +J(1,1)\nu^-\otimes\nu^+.
\end{equation}
For a bounded test $F$ define $\mathsf D_{\rm test}F$ by subtracting its two
integrated marginals and adding its product mean.  Then
\[
 \int F\,d(\mathsf D_{\rm cur}J)=
 \int \mathsf D_{\rm test}F\,dJ.
\]
Thus $J^\circ$ is double centered and the current projection is the adjoint,
not the same operator, as test centering.  If the complete response current
already has zero marginals, then
\begin{equation}
\label{eq:global-recentering-sum}
 J^\circ=J.
\end{equation}
Only after this identity is formed in physical coordinates is the singular
part refined by face labels.  If the original physical amplitudes obey
\eqref{eq:global-ledger-sum} and the references have fixed $C^\alpha$
densities, then
\begin{equation}
\label{eq:recenter-norm-ledger}
 \|\mathsf D_{\rm cur}J\|_{\rm TV}
 \le4\|J\|_{\rm TV},\qquad
 \sum_\lambda\|J_\lambda^\circ\|_{\rm phys,R}<\infty .
\end{equation}
\end{lemma}

\begin{proof}
Fubini proves the adjoint identity and each marginal has total variation at
most $\|J\|_{\rm TV}$.  This gives the factor four.  Marginalization and
tensoring with a fixed smooth probability preserve every scaled mark in the
physical ledger.  Absolute summation of \eqref{eq:global-ledger-sum} proves
the second assertion.  No Banach space of graph currents or projective tensor
completion is introduced.
\end{proof}

\begin{definition}[Branch-section interface $\mathrm{BR}_{\rm sec}$]
\label{def:branch-section-interface}
On every clean graph rectangle, the dominant endpoint coordinates admit
branchwise sections
\[
 \iota^-_\lambda:X\mapsto(X,\zeta^-_\lambda(X)),\qquad
 \iota^+_\lambda:Y\mapsto(Y,\zeta^+_\lambda(Y))
\]
inside the same homogeneous inverse branches as the physical graph trace.
Their images, in the appropriate forward and reversed orientations, are
homogeneous measured standard curves.  Their arclength/coordinate densities,
endpoint $Z$ values, and scaled H\"older marks are bounded by
$\operatorname{Mark}(\lambda)\{1+Z(\mathcal G^-_\lambda)+
Z(\mathcal G^+_\lambda)\}$.  For $h_M=Q^Mg$ and
$f_N=(Q^*)^Nf$, branchwise contraction gives
\begin{align}
 \int|h_M(x_\lambda(\bar s))-
 h_M(\iota^-_\lambda(X_\lambda(\bar s)))|\,d|\omega_\lambda|
 &\le Ce^{-c_-M}\|g\|_{C^1}|\omega_\lambda|,\notag\\
 \int|f_N(y_\lambda(\bar s))-
 f_N(\iota^+_\lambda(Y_\lambda(\bar s)))|\,d|\omega_\lambda|
 &\le Ce^{-c_+N}\|f\|_{C^1}|\omega_\lambda|,
 \label{eq:branch-section-contraction}
\end{align}
where $d\omega_\lambda=M_\lambda w_\lambda d\bar s$.
This is a same-branch condition: no conditional measure is averaged across
different symbolic pasts or singularity cuts.
\end{definition}

\begin{lemma}[Dominant-coordinate branch-section reduction]
\label{lem:dominant-coordinate-reduction}
Assume $\mathrm{BR}_{\rm sec}$.  After assembly and the recentering of Lemma
\ref{lem:global-label-recentering}, retain the rank-one and regular-product
blocks as ordinary product kernels and evolve both factors by their centered
operators.
On a clean graph rectangle in the diagonal band
$D^{-1}M<N<AM$, the graph-current pairing decomposes as
\begin{equation}
\label{eq:dominant-coordinate-decomposition}
 B_{M,N}=B^{\rm dom,dom}_{M,N}
 +B^{\rm dom,rem}_{M,N}+B^{\rm rem,dom}_{M,N}
 +B^{\rm rem,rem}_{M,N}.
\end{equation}
The first term is the scalar quotient integral obtained by evaluating the two
tests on $\iota^-_\lambda(X)$ and $\iota^+_\lambda(Y)$.
Every other term is bounded by
\begin{equation}
\label{eq:contracted-coordinate-error}
 Ce^{-c\max\{M,N\}}\|g\|_{C^1}\|f\|_{C^1}.
\end{equation}
The global joint constant and axis blocks vanish or remain in the explicitly
separated rank-one summands before the graph-rectangle sum is estimated.
Every graph-axis mode created by a boundary cutoff is uniformly
nonstationary and is charged to its cutoff norm.
\end{lemma}

\begin{proof}
In stable--unstable endpoint coordinates, use the same-branch sections of
$\mathrm{BR}_{\rm sec}$ and telescope the two replacements.
Equation~\eqref{eq:branch-section-contraction} gives remainders bounded by
$C\lambda_-^M\|g\|_{C^1}$ and $C\lambda_+^N\|f\|_{C^1}$ after the physical
amplitude is integrated.  In the diagonal band these imply
\eqref{eq:contracted-coordinate-error}.  The two frozen functions depend
only on $X,Y$, proving \eqref{eq:dominant-coordinate-decomposition}.  Global
projection separates the constant and two axis blocks before any
label-dependent chart is chosen.  They are estimated by their explicit
product formulas, not redescribed as graph currents; Lemma
\ref{lem:rank-one-product} supplies product decay.  A graph boundary
cutoff can reintroduce an axis term, but it has phase $\xi_1X$ or $\xi_2Y$;
\eqref{eq:dominant-coordinate-normalization} gives the physical derivative
lower bound, and its cutoff ledger is absolutely summable by
\eqref{eq:recenter-norm-ledger}.
\end{proof}

\begin{proposition}[Scaled physical face density]
\label{prop:scaled-face-density}
On a homogeneous physical face cell
\(Q=t_Q+\ell_Q[-1,1]\), the collision--SRB conditional is
\(d\mu_Q^u=\rho_Q(t)\,dt\), where, before normalization,
\begin{equation}
\label{eq:face-conditional-density}
 \rho_Q^{\rm raw}(t)=c_\mu\cos\varphi(t)\,J_{\rm coa}(t)
 J_{\rm hol}(t)J_{\rm chart}(t),
 \qquad J_{\rm coa}=|\nabla G|^{-1}.
\end{equation}
Put
\(\widehat\rho_Q(u)=\ell_Q\rho_Q(t_Q+\ell_Qu)\).  There is a mark
\begin{equation}
\label{eq:density-mark-definition}
 \operatorname{Mark}(Q)=
 1+\mathfrak h_Q+\mathfrak b_Q+
 [\log J_{\rm hol}]_{\alpha,Q}
 +[\log J_{\rm coa}]_{\alpha,Q}
 +[\log J_{\rm chart}]_{\alpha,Q}
 +\|w_Q\|_{\mathcal A^\alpha}
\end{equation}
such that
\begin{equation}
\label{eq:scaled-density-bound}
 C^{-1}\operatorname{Mark}(Q)^{-1}
 \le\widehat\rho_Q\le C\operatorname{Mark}(Q),\qquad
 [\log\widehat\rho_Q]_\alpha
 \le C\operatorname{Mark}(Q).
\end{equation}
For some \(p_d>2\),
\begin{equation}
\label{eq:density-mark-moment}
 \sum_Q M_Q\operatorname{Mark}(Q)^{p_d}<\infty
\end{equation}
uniformly in the small table family.  The same statement holds after the
positive/negative decomposition of a signed face amplitude.
\end{proposition}

\begin{proof}
In collision coordinates the invariant density is
\(c_\mu\cos\varphi\,dr\,d\varphi\).  Coarea along the physical graph, followed
by stable projection to the gate coordinate, gives
\eqref{eq:face-conditional-density}.  On a homogeneous strip,
\(\cos\varphi\), the chart Jacobian, and the holonomy Jacobian have the
standard scaled log-Hölder distortion.  Dividing by
\(\int_Q\rho_Q^{\rm raw}\) gives \eqref{eq:scaled-density-bound}; the factor
\(\ell_Q\) removes the cell scale.

The homogeneity contribution has a fixed polynomial cost, while the strip
mass is exponentially summable.  The holonomy and chart contributions are
summed by the context moment, and the amplitude contribution by the marked
face ledger.  Generalized Hölder with slightly smaller exponents proves
\eqref{eq:density-mark-moment}.  For a signed amplitude \(q\), write
\(q=(2\|q\|_\infty+q)-2\|q\|_\infty\); multiplication by the positive masses
restores the scaled norm and introduces no reciprocal small-mass factor.
\end{proof}

\begin{lemma}[Fractional nonstationary phase]
\label{lem:fractional-nonstationary-phase}
Let \(I\) be an interval, \(a\in C^\alpha(I)\), and
\(\phi\in C^{1+\alpha}(I)\).  On each component of the support of \(a\),
suppose \(\phi'\) has one sign and \(|\phi'|\ge\lambda>0\).  Put
\(\Phi=\sum_C\int_C|\phi'(t)|\,dt\), the total phase length of the support.
For $C=(l_C,r_C)$ define the endpoint functional
$\mathfrak E(a)=\sum_C(|a(l_C+)|+|a(r_C-)|)$.
\begin{equation}
\label{eq:fractional-nsp}
 \left|\int_Ie^{i\phi(t)}a(t)\,dt\right|
 \le C_\alpha\left\{
 \frac{\mathfrak E(a)}{\lambda}
 +\frac{\Phi[a]_\alpha}{\lambda^{1+\alpha}}
 +\frac{\Phi\|a\|_\infty[\phi']_\alpha}
             {\lambda^{2+\alpha}}\right\}.
\end{equation}
The estimate adds over a countable union of sign components without a
component-count loss.
\end{lemma}

\begin{proof}
On one sign component use \(y=\phi(t)\) and write
\(b(y)=a(\phi^{-1}y)/\phi'(\phi^{-1}y)\).  Then
\[
 \|b\|_\infty\le\|a\|_\infty/\lambda,\qquad
 [b]_\alpha\le
 [a]_\alpha\lambda^{-1-\alpha}
 +\|a\|_\infty[\phi']_\alpha\lambda^{-2-\alpha}.
\]
For a shift \(h=\pi\), the identity
\[
 2\int e^{iy}b(y)\,dy
 =\int e^{iy}\{b(y)-b(y+h)\}\,dy+\text{endpoint terms}
\]
gives $C\Phi[b]_\alpha$ plus the two endpoint values of $b$ on that
component.  Since $|b|\le |a|/\lambda$, summing the endpoint values gives
$\mathfrak E(a)/\lambda$, while the sum of the component $y$-lengths is
$\Phi$.  This proves \eqref{eq:fractional-nsp}; no number-of-components
estimate is hidden in the supremum term.
\end{proof}

\begin{lemma}[Rank-one product under double centering]
\label{lem:rank-one-product}
If \(\eta_\pm\) are admissible dual collision currents and
\(R=\eta_-\otimes\eta_+\), then
\[
 R(Q_-^mg\otimes Q_+^nf)
 =\eta_-(Q_-^mg)\eta_+(Q_+^nf).
\]
Thus the two one-sided gaps give \(C\tau^m\vartheta^n\).
\end{lemma}

\begin{proof}
This is the universal property of the projective tensor product followed by
the two admissible dual bounds.  In the principal proof the rank-one
corrections are instead included in the low-frequency functional of
Proposition~\ref{prop:functional-correlation-low-frequency}; this lemma is
only an algebraic consistency check.
\end{proof}

\begin{lemma}[Balanced--unbalanced exhaustion in physical frequency]
\label{lem:balanced-unbalanced-exhaustion}
Assume the clean-gate chart bounds
\begin{equation}
\label{eq:primitive-dominant-chart-bounds}
 c_0a_X\le |X'|\le C_0a_X,\qquad
 c_0a_Y\le |Y'|\le C_0a_Y
\end{equation}
with the fixed orientation signs on one homogeneous face cell.  Put
$C_{\rm bal}=4C_0/c_0$.  Every nonzero scalar mode belongs to exactly one of
the following regions (with the boundary assigned once):
\begin{enumerate}[label=\textup{(\roman*)}]
\item $C_{\rm bal}^{-1}\le
 |\xi_1|a_X/(|\xi_2|a_Y)\le C_{\rm bal}$ (balanced);
\item $|\xi_1|a_X>C_{\rm bal}|\xi_2|a_Y$ (past-dominant);
\item $|\xi_2|a_Y>C_{\rm bal}|\xi_1|a_X$ (future-dominant).
\end{enumerate}
On either unbalanced region the phase
$\psi_\xi=\xi_1X+\xi_2Y$ has one-sign derivative after at most the fixed
orientation cut and
\begin{equation}
\label{eq:unbalanced-derivative-lower-bound}
 |\psi_\xi'|\ge c\max\{|\xi_1|a_X,|\xi_2|a_Y\}.
\end{equation}
Consequently the complete unbalanced contribution is bounded by
\begin{equation}
\label{eq:unbalanced-mode-saving}
 C\operatorname{Mark}(Q)^2\{1+Z(\mathcal G_Q)\}
 \max\{|\xi_1|a_X,|\xi_2|a_Y\}^{-\alpha},
\end{equation}
and its two-dimensional lattice sum converges under the same observable
regularity used for the balanced sum.
\end{lemma}

\begin{proof}
Consider the past-dominant case.  The orientation table and
\eqref{eq:primitive-dominant-chart-bounds} give
\[
 |\psi_\xi'|
 \ge c_0|\xi_1|a_X-C_0|\xi_2|a_Y
 \ge \tfrac12c_0|\xi_1|a_X.
\]
The other case is symmetric.  Cut only at the fixed orientation endpoints;
no frequency-dependent sign partition is introduced.  Apply Lemma
\ref{lem:fractional-nonstationary-phase} with this lower bound and with the
actual amplitude and endpoint ledger.  Its three terms are bounded by the
right side of \eqref{eq:unbalanced-mode-saving}; the exponent $\alpha$ may be
decreased once to absorb fixed chart losses.  Multiplication by the Fourier
coefficients of the two $C^{r_{\rm obs}}$ tests and summation over the two
unbalanced cones is finite as soon as $r_{\rm obs}>2$; the stronger balanced
threshold imposed below therefore suffices.  The three regions are disjoint
and exhaustive.
\end{proof}

\begin{proposition}[Central balanced estimate with explicit budget]
\label{prop:central-phase-verification}
Let \(\Theta\) be the energy exponent and set
\[
 A_{\rm frac}=3+\alpha,\qquad
 0<\delta<\frac{\alpha}{2A_{\rm frac}},\qquad
 \beta_{\rm cen}=
 \min\{\delta\Theta/2,\alpha-A_{\rm frac}\delta\}>0.
\]
Then every balanced nonzero graph mode has saving
\(C\Xi^{-\beta_{\rm cen}}\) times a fixed power of
\(\operatorname{Mark}(Q)\).  If the observable regularity satisfies
\(r_{\rm obs}>2+\beta_{\rm cen}\), the complete lattice sum is exponentially
small in the two induced depths.  Identity
\eqref{eq:section-quotient-low-high} assigns every localized quotient mode
exactly once.  Low graph modes, rank-one corrections, and
the regular product block are treated together by
Proposition~\ref{prop:functional-correlation-low-frequency}.
\end{proposition}

\begin{proof}
Take \(r=\Xi^{-\delta}\).  The energy theorem bounds the resonant mass by
\(C\Xi^{-\delta\Theta/2}+Ce^{-cN}\).  On its complement use two smooth
cutoffs in
\(\psi_\xi'/(\Xi r)\), one for each sign.  They vanish on the resonant
boundaries.  Sum adjacent physical cells before taking absolute values, so
their artificial endpoint terms cancel with opposite orientation.  The
remaining face, singularity, and homogeneity endpoints obey
\begin{equation}
\label{eq:phase-endpoint-ledger}
 \sum_Q\mathfrak E(a_Q)
 \le Cr^{-1}\sum_QM_Q\operatorname{Mark}(Q)
       \{1+Z(\mathcal G_Q)\},
\end{equation}
by Proposition~\ref{prop:radial-direct-Z}.  Proposition
\ref{prop:scaled-face-density} gives
\[
 \|a_\pm\|_\infty+[a_\pm]_\alpha
 \le C\operatorname{Mark}(Q) r^{-1}.
\]
Furthermore
\[
 \lambda=c\Xi r,\qquad
 \Phi\le C\Xi,\qquad
 [\psi_\xi']_\alpha\le C\Xi\operatorname{Mark}(Q).
\]
Lemma~\ref{lem:fractional-nonstationary-phase} therefore gives
\begin{align*}
 |\mathcal I_Q(\xi)|
 &\le C\operatorname{Mark}(Q)^2
 \{(1+Z(\mathcal G_Q))\Xi^{-1}r^{-2}
        +\Xi^{-\alpha}r^{-(3+\alpha)}\}\\
 &\le C\operatorname{Mark}(Q)^2(1+Z(\mathcal G_Q))
       \Xi^{-\beta_{\rm cen}}.
\end{align*}
The joint estimate \eqref{eq:joint-mark-Z-ledger} sums the cells without
separating the density and boundary moments.

For completeness, put $K=e^{\eta N}$ with
$0<\eta<\min\{c_-,c_+\}/(2A_{\rm fr})$.  After global centering and the
dominant-coordinate reduction, the face-adapted projector
$\mathsf S_K^{\rm dom}$ retains precisely the scalar modes satisfying
$|\xi_1|a_X\le K$ and $|\xi_2|a_Y\le K$.  Proposition
\ref{prop:functional-correlation-low-frequency} treats this finite sum with
finite-rank loss $K^{A_{\rm fr}}$.  Every balanced mode of the complementary dominant
sum has $\Xi>K$ by definition, so the preceding saving is
$e^{-\eta\beta_{\rm cen}N}$.  Contracted-coordinate terms were removed before
the cutoff and are exponentially small by Lemma
\ref{lem:dominant-coordinate-reduction}.  The difference between the graph
trace and its same-branch section is already part of that contracted error;
no fixed-fiber average is used.  The two unbalanced scalar cones are
estimated, with their full lattice sum, by Lemma
\ref{lem:balanced-unbalanced-exhaustion}.  The two-dimensional dominant lattice sum converges because
\(r_{\rm obs}>2+\beta_{\rm cen}\).  Lemma~\ref{lem:gate-buffer} supplies the
resolving depth.  Global marginal corrections are inserted through the typed
current/test adjoint identity before any absolute value is taken.  The
regular two-dimensional density is not placed in the face lattice:
Corollary~\ref{cor:regular-functional-correlation} uses its independent
$W^{R,1}$ smoothing tail.
\end{proof}

\subsection{Low frequencies after dominant-coordinate reduction}

The cutoff is imposed only after the graph term has been reduced to its two
scalar dominant coordinates.  This avoids the false implication that a large
ambient Fourier vector must have a large restriction to a one-dimensional
face.  Contracted coordinates are estimated before the cutoff.

\begin{definition}[Dominant event norm and separate dynamic-H\"older norm]
\label{def:event-current-norm}
Fix a finite physical product atlas and assemble coincident artificial faces
before inserting its partition of unity.  Let $\nu^\pm$ be the two
transported stationary smooth references and set
$dm_{M\times M}=d\nu^-\otimes d\nu^+$.  For a face label $\lambda$ let
$a_{X,\lambda},a_{Y,\lambda}>0$ be the two scales in
\eqref{eq:dominant-coordinate-normalization} and put
\[
 \mathfrak d_\lambda=1+a_{X,\lambda}^{-1}+a_{Y,\lambda}^{-1}.
\]
Write
\[
 \mathcal K_{\rm ev}=k_{\rm reg}\,dm_{M\times M}
 +\sum_\lambda M_\lambda(H_\lambda)_*(w_\lambda\,d\bar s)
\]
and let $q_{\rm sec}\ge1$ be the fixed polynomial recovery exponent for one
oscillatory section current in
Lemma~\ref{lem:oscillatory-section-recovery}.  Put
$A_{\rm sec}:=2+2q_{\rm sec}\ge4$.
For an integer $1\le R\le3$ allowed by the $C^3$ table regularity, define
\begin{equation}
\label{eq:event-physical-norm}
 \|\mathcal K_{\rm ev}\|_{\rm phys,R}:=
 \sum_\chi\|\zeta_\chi k_{\rm reg}\|_{W^{R,1}}
 +\sum_\lambda M_\lambda\|w_\lambda\|_{\mathcal A^\alpha}
 \mathfrak d_\lambda^{A_{\rm sec}}\operatorname{Mark}(\lambda)^2
 \{1+Z(\mathcal G_\lambda^-)+Z(\mathcal G_\lambda^+)\}.
\end{equation}
For bounded $F:M^q\to\mathbb C$ let
\begin{equation}
\label{eq:separate-dynamic-holder-norm}
 \|F\|_{\rm SH}:=\|F\|_\infty+
 \sum_{i=1}^q\sup_{x_j:j\ne i}\sup_{x_i\ne y_i}
 \frac{|F(\ldots,x_i,\ldots)-F(\ldots,y_i,\ldots)|}
      {\theta_0^{s_i(x_i,y_i)}}.
\end{equation}
The stable or unstable separation time appropriate to the $i$th slot is used,
as in \cite[Definition~2.2]{LS}.
\end{definition}

We now define the low-frequency object on the scalar quotient on which it
actually lives.  Fix an even cutoff $\chi\in C_c^\infty((-2,2))$ with
$\chi=1$ on $[-1,1]$.  After the branch-section reduction, the assembled
graph piece induces the finite signed scalar measure
\begin{equation}
\label{eq:dominant-pushforward-measure}
 \nu_\lambda=(X_\lambda,Y_\lambda)_*
       (M_\lambda w_\lambda\,d\bar s),\qquad
 \widehat\nu_\lambda(k,l)=
 \int e^{-i(kX_\lambda+lY_\lambda)}M_\lambda w_\lambda\,d\bar s .
\end{equation}
Define the anisotropic multiplier and its polynomial density by
\begin{align}
 m_{\lambda,K}(k,l)
 &=\chi(a_{X,\lambda}k/K)\chi(a_{Y,\lambda}l/K),
 \label{eq:dominant-physical-cutoff}\\
 A_{\lambda,K}^{\rm dom}(x,y)
 &=(2\pi)^{-2}\sum_{(k,l)\in\mathbb Z^2}
 m_{\lambda,K}(k,l)\widehat\nu_\lambda(k,l)e^{i(kx+ly)}.
 \label{eq:canonical-dominant-polynomial}
\end{align}
This is a density on the scalar quotient and nowhere else.  Define the two
mode currents on the branch sections by
\begin{align}
 \eta^-_{\lambda,k}(h)
 &:=\int h(\iota^-_\lambda(X))\chi^-_\lambda(X)e^{ikX}\,dX,
 \label{eq:past-section-mode-current}\\
 \eta^+_{\lambda,l}(f)
 &:=\int f(\iota^+_\lambda(Y))\chi^+_\lambda(Y)e^{ilY}\,dY,
 \label{eq:future-section-mode-current}
\end{align}
where the fixed scalar cutoffs are part of the localized quotient chart.
They equal one on a neighborhood of
$\operatorname{supp}\nu_\lambda$ and their supports stay inside the
homogeneous section domains.
The branch-section lift $\mathfrak L^{\rm sec}_{\lambda,K}$ is
\begin{equation}
\label{eq:section-lifted-low-block}
 \int h(x)f(y)\,d\mathfrak L^{\rm sec}_{\lambda,K}(x,y)
 :=\int h(\iota^-_\lambda(X))f(\iota^+_\lambda(Y))
       \chi^-_\lambda(X)\chi^+_\lambda(Y)
       A_{\lambda,K}^{\rm dom}(X,Y)\,dX\,dY.
\end{equation}
If $\mathsf P_{\lambda,K}$ denotes the even quotient multiplier in
\eqref{eq:dominant-physical-cutoff} and
$H(X,Y)=h(\iota^-_\lambda(X))f(\iota^+_\lambda(Y))$, then self-adjointness of
the multiplier and the support clause give the exact localized identity
\begin{equation}
\label{eq:section-quotient-low-high}
 \int H\,d\nu_\lambda-
 \int \chi^-_\lambda\chi^+_\lambda H A_{\lambda,K}^{\rm dom}\,dX\,dY
 =\int (I-\mathsf P_{\lambda,K})
       (\chi^-_\lambda\chi^+_\lambda H)\,d\nu_\lambda.
\end{equation}
Thus the only localization error is the same quotient high-pass term treated
by the balanced/unbalanced phase estimates; it is not an unrecorded leakage
of the Fourier polynomial outside the branch chart.
Fourier inversion on the quotient gives the exact
finite-rank identity
\begin{equation}
\label{eq:section-current-finite-rank}
 \mathfrak L^{\rm sec}_{\lambda,K}
 =(2\pi)^{-2}\sum_{k,l}
 m_{\lambda,K}(k,l)\widehat\nu_\lambda(k,l)
 \eta^-_{\lambda,k}\otimes\eta^+_{\lambda,l}.
\end{equation}
Thus finite rank is obtained in the tensor product of two concrete measured
standard-curve currents.  The original graph current is not claimed to be
nuclear or absolutely continuous.

For a finite signed kernel $\sigma$ on $M\times M$ put
\begin{equation}
\label{eq:kernel-operator-identification}
 \langle\operatorname{Op}(\sigma)h,f\rangle
 :=\int_{M\times M}h(x)f(y)\,d\sigma(x,y).
\end{equation}
\begin{lemma}[Oscillatory section currents are recoverable]
\label{lem:oscillatory-section-recovery}
For each mode current in
\eqref{eq:past-section-mode-current}--
\eqref{eq:future-section-mode-current}, its real and imaginary parts are
finite signed differences of positive measured standard families.  Their
total masses are bounded by the section mark, their scaled density marks are
at most
$CK\mathfrak d_\lambda\operatorname{Mark}(\lambda)$, and they are proper after
\begin{equation}
\label{eq:section-mode-recovery-time}
 R_{\lambda,K}^{\pm}\le C_0+C_1\log(K\mathfrak d_\lambda)
 +C_2\log\{1+Z(\mathcal G_\lambda^\pm)\}.
\end{equation}
The same assertion holds for the reversed collision map.
Moreover, if $\|\cdot\|_{\rm SP,\pm}$ is the centered one-sided
standard-family mixing dual norm including the finite recovery window, then
\begin{equation}
\label{eq:section-mode-dual-cost}
 \|\eta^\pm_{\lambda,k}\|_{\rm SP,\pm}
 \le C(K\mathfrak d_\lambda)^{q_{\rm sec}}
 \operatorname{Mark}(\lambda)\{1+Z(\mathcal G_\lambda^\pm)\}.
\end{equation}
\end{lemma}

\begin{proof}
Relative to the positive coordinate/arclength density on the section, the
multiplier is a fixed cutoff times $e^{ikX}$.  Split real and imaginary parts
and add twice their supremum to obtain positive densities, exactly as in
Lemma~\ref{lem:signed-observable-multiplier}.  One scalar derivative costs
$C|k|\le CK/a_{X,\lambda}$ (and similarly in $Y$), which gives the displayed
scaled mark.  The density recursion contracts this mark in
$O(\log(K\mathfrak d_\lambda))$ iterates, while the boundary recursion costs
$O(\log(1+Z))$.  Forward/reverse proper-family growth and coupling are uniform
on the compact table class, proving \eqref{eq:section-mode-recovery-time}.
Exponential convergence delayed by that time costs a fixed power of
$K\mathfrak d_\lambda$; enlarge it once to cover the finite recovery window.
This defines $q_{\rm sec}$ and proves \eqref{eq:section-mode-dual-cost}.
\end{proof}

There are two different centering operations, and we keep their types
separate.  If $\sigma$ is a finite signed measure on $M\times M$, write
$\sigma_\pm=(\operatorname{pr}_\pm)_*\sigma$ and
$|\sigma|_0=\sigma(M\times M)$, and define
\begin{equation}
\label{eq:current-centering}
 \mathsf D_{\rm cur}\sigma
 :=\sigma-\sigma_-\otimes\nu^+-\nu^-\otimes\sigma_+
       +|\sigma|_0\nu^-\otimes\nu^+.
\end{equation}
Its test-side adjoint is
\begin{align}
 \mathsf D_{\rm test}F(x,y)
 :=F(x,y)&-\int F(x,y')\,d\nu^+(y')
          -\int F(x',y)\,d\nu^-(x') \notag\\
 &+\int F(x',y')\,d\nu^-(x')d\nu^+(y'),
 \label{eq:test-centering}
\end{align}
so that
\begin{equation}
 \label{eq:centering-adjoint-identity}
 \int F\,d(\mathsf D_{\rm cur}\sigma)
 =\int \mathsf D_{\rm test}F\,d\sigma.
\end{equation}
In particular a current is never silently fed to an operator defined on
functions.  On the regular density use the ordinary
four-dimensional convolution $\rho_K^{(4)}*k_{\rm reg}$ and define
\begin{align}
\label{eq:dominant-centered-smoother}
\mathsf S_K^{\rm dom}\mathcal K_{\rm ev}:=
 \mathsf D_{\rm cur}\left(
 [\rho_K^{(4)}*k_{\rm reg}]dm_{M\times M}
 +\sum_\lambda\mathfrak L^{\rm sec}_{\lambda,K}\right),
 \notag\\
 \mathsf H_K^{\rm dom}(\mathsf D_{\rm cur}\mathcal K_{\rm ev}):=
 \mathsf D_{\rm cur}\mathcal K_{\rm ev}-
 \mathsf S_K^{\rm dom}\mathcal K_{\rm ev}.
\end{align}
All artificial endpoints are assembled before this definition.  Expanding
$\mathsf D_{\rm cur}$ centers each pair of section-mode currents and includes
the two marginal and constant product blocks explicitly.
The definition of $\mathsf H_K^{\rm dom}$ gives the exact operator decomposition
\begin{equation}
\label{eq:current-operator-low-high-identity}
 \operatorname{Op}(\mathsf D_{\rm cur}\mathcal K_{\rm ev})
 =\operatorname{Op}(\mathsf S_K^{\rm dom}\mathcal K_{\rm ev})
  +\operatorname{Op}\!\left(
       \mathsf H_K^{\rm dom}(\mathsf D_{\rm cur}\mathcal K_{\rm ev})\right).
\end{equation}

\begin{theorem}[Branch-section low projection and finite rank]
\label{thm:canonical-dominant-smoothing}
Put $A_{\rm fr}=\max\{6,A_{\rm sec}\}$.  For $K\ge2$ the centered current in
\eqref{eq:dominant-centered-smoother} has an exact finite-rank expansion
\begin{equation}
\label{eq:finite-rank-low-kernel}
 \mathsf S_K^{\rm dom}\mathcal K_{\rm ev}
 =\sum_{j\in\mathcal J_K}c_j
       (\eta_j^-)^\circ\otimes(\eta_j^+)^\circ,
\end{equation}
where $\eta^\circ=\eta-\eta(1)\nu^\pm$.  Face factors are the concrete
section currents \eqref{eq:past-section-mode-current}--
\eqref{eq:future-section-mode-current}; regular-density factors are ordinary
absolutely continuous currents.  With the one-sided mixing norms of
\eqref{eq:section-mode-dual-cost},
\begin{equation}
\label{eq:finite-rank-projective-cost}
 \sum_{j\in\mathcal J_K}|c_j|
 \|(\eta_j^-)^\circ\|_{\rm SP,-}
 \|(\eta_j^+)^\circ\|_{\rm SP,+}
 \le CK^{A_{\rm fr}}\|\mathcal K_{\rm ev}\|_{\rm phys,R}.
\end{equation}
The constant is independent of return and grazing truncations, and
\eqref{eq:current-operator-low-high-identity} holds before any absolute value
is taken.  This is a finite-rank statement only after scalar cutoff; no
nuclear or absolutely continuous representation of the original graph current
is asserted.
\end{theorem}

\begin{proof}
The coefficient bound
$|\widehat\nu_\lambda(k,l)|\le
M_\lambda\|w_\lambda\|_{L^1}$ follows directly from
\eqref{eq:dominant-pushforward-measure}.  The multiplier has at most
$CK^2\mathfrak d_\lambda^2$ nonzero pairs.  Formula
\eqref{eq:section-current-finite-rank} gives their exact product-current
expansion.  By \eqref{eq:section-mode-dual-cost}, the product of the two
one-sided norms costs at most
$C(K\mathfrak d_\lambda)^{2q_{\rm sec}}$ times the section mark--$Z$ factor.
Thus the face block costs at most
$CK^{A_{\rm sec}}\mathfrak d_\lambda^{A_{\rm sec}}$ times the corresponding
summand of \eqref{eq:event-physical-norm}.  The regular
four-dimensional convolution has $O(K^4)$ modes and the two factors cost at
most $K^2$, giving exponent $6$.  Taking the maximum with $A_{\rm sec}$ and
absolute
summation in \eqref{eq:event-physical-norm} proves
\eqref{eq:finite-rank-projective-cost}.  Formula
\eqref{eq:current-centering} replaces each factor by its centered version,
and \eqref{eq:centering-adjoint-identity} proves the four marginal identities.
The definition of the remainder proves the exact low/high operator identity;
Fourier inversion is used only on the scalar quotient.  The preliminary
replacement of the physical graph trace by the branch sections is charged
separately to \eqref{eq:contracted-coordinate-error} before the cutoff.
\end{proof}

\begin{lemma}[Regular high-frequency tail]
\label{lem:regular-high-tail}
For the regular two-dimensional density,
\begin{equation}
\label{eq:regular-high-tail}
 \|(I-\rho_K^{(4)}*)k_{\rm reg}\|_{L^1(M\times M)}
 \le CK^{-R}\|k_{\rm reg}\|_{W^{R,1}(M\times M)}.
\end{equation}
Consequently its pairing with $Q^mg$ and $Q^nf$ is bounded by the right-hand
side times $\|g\|_\infty\|f\|_\infty$, uniformly in $m,n$.
\end{lemma}

\begin{proof}
This is the standard approximation-of-the-identity estimate in $L^1$; Markov
normalization gives $\|Q^mh\|_\infty\le2\|h\|_\infty$.
\end{proof}

\begin{proposition}[Central low-frequency estimate by two-sided mixing]
\label{prop:functional-correlation-low-frequency}
Let $c_-,c_+>0$ be the uniform past and future two-state mixing rates supplied
by the finite-horizon functional-correlation hypothesis.  In the diagonal band put
$N=\max\{m,n\}$ and $K=e^{\eta N}$.  Then
\begin{equation}
\label{eq:functional-correlation-low-frequency}
 \left|\left\langle Q^n
 \operatorname{Op}(\mathsf S_K^{\rm dom}\mathcal K_{\rm ev})Q^mg,
 f\right\rangle\right|
 \le C e^{-c_-m-c_+n+A_{\rm fr}\eta N}
 \|g\|_{C^r}\|f\|_{C^r},
\end{equation}
provided $0<\eta<\min\{c_-,c_+\}/(2A_{\rm fr})$.  In particular, since
$N\le m+n$, the right side is bounded by $Ce^{-c_{\rm low}(m+n)}$ for
$c_{\rm low}=\min\{c_-,c_+\}-A_{\rm fr}\eta>0$.
\end{proposition}

\begin{proof}
Insert \eqref{eq:finite-rank-low-kernel} into the product-current pairing.  By
Fubini and \eqref{eq:kernel-operator-identification},
\[
 \left\langle Q^n\operatorname{Op}(\mathsf S_K^{\rm dom}\mathcal K_{\rm ev})
 Q^mg,f\right\rangle
 =\sum_jc_j(\eta_j^-)^\circ(Q^mg)
              (\eta_j^+)^\circ(Q^{*n}f).
\]
These are two independent one-time standard-family correlations; there is no
condition $y=Tx$.  Forward and reversed proper-family coupling, including
the recovery window in \eqref{eq:section-mode-dual-cost}, bounds the factors by
$Ce^{-c_-m}\|g\|_{C^r}\|(\eta_j^-)^\circ\|_{\rm SP,-}$ and
$Ce^{-c_+n}\|f\|_{C^r}\|(\eta_j^+)^\circ\|_{\rm SP,+}$.  Sum with
\eqref{eq:finite-rank-projective-cost} and use $K=e^{\eta N}$.
\end{proof}

\begin{corollary}[Regular defect via functional correlation]
\label{cor:regular-functional-correlation}
The regular part of $K_V$ has central-band exponential decay: its finite-rank
low part is covered by Proposition~\ref{prop:functional-correlation-low-frequency}, and
Lemma~\ref{lem:regular-high-tail} gives $K^{-R}=e^{-\eta RN}$ for its high
part.  No tensor, wavelet, or face-phase estimate is used.
\end{corollary}

\subsection{Radial faces: direct flux boundary control}

We retain every connected component.  Let
$\mathcal W_\lambda(j,k)=\{W_{\lambda,j,k,r}\}_r$ be the maximal components
of one coarea face with incoming strip $j$ and outgoing strip $k$.  Put
\[
 p_{\lambda,j,k,r}=\int_{W_{\lambda,j,k,r}}\rho^{\rm raw}_\lambda\,ds,
 \quad M_\lambda=\sum_{j,k,r}p_{\lambda,j,k,r},\quad
 \bar p_{\lambda,j,k,r}=p_{\lambda,j,k,r}/M_\lambda.
\]
Thus $\bar p$, not $p$, is the weight of the normalized standard family.

\begin{lemma}[Oriented refinement bookkeeping]
\label{lem:cone-monotone-components}
Assume a face has been cut into $N_{\rm or}(\lambda)$ signed charts on each
of which either the forward or the reversed input and output angle
coordinates are monotone.  Then for fixed incoming/outgoing signed
homogeneity ranks $(j,k)$ there are at most $4N_{\rm or}(\lambda)$ connected
components and
\begin{equation}
\label{eq:derived-weighted-component-complexity}
 \sum_{j,k,r}\frac1{(1+j)^2(1+k)^2}
 \le C N_{\rm or}(\lambda).
\end{equation}
The existence of such a cover and every estimate for $N_{\rm or}$ are not
deduced from cone transversality or a qualitative jet condition; their
weighted form is an explicit part of $\mathrm{FZ}_{\rm geom}$ below.
\end{lemma}

\begin{proof}
On an unstable graph the cone bounds give $dr/dt$ and $d\varphi/dt$ one fixed
sign after the orientation of $t$ is chosen.  Invariance of the unstable cone
gives the same statement after one collision; the reversed cone gives the
backward version.  A signed homogeneity strip is an interval in the angle
coordinate.  The preimages of an incoming and an outgoing signed strip are
therefore intervals on the oriented chart, and their intersection is an
interval.  The two grazing signs in each time orientation account for the
factor four.  Summing over charts and then over $j,k$ proves
\eqref{eq:derived-weighted-component-complexity}.  This proves only the
bookkeeping implication from an already constructed oriented cover; it makes
no analytic zero-count assertion.
\end{proof}

\begin{lemma}[Full tangent coarea Jacobian]
\label{lem:two-flux-coarea}
Let $\gamma(t)=(r_-(t),\varphi_-(t))$ be a unit-speed parametrization of an
input coarea face and let $W=T\gamma$.  On every oriented component the
unnormalized density with respect to output arclength is
\begin{equation}
\label{eq:radial-two-flux-density}
 \rho^{\rm raw}_{W,\lambda}(T\gamma(t))=
 \frac{|a_\lambda|J_{{\rm chart},\lambda}J_{{\rm hol},\lambda}
       \cos\varphi_-(t)}
 {|\nabla G_\lambda(\gamma(t))|\,
       |D T(\gamma(t))\gamma'(t)|}.
\end{equation}
If output impact coordinate $r_+$ is used instead of arclength, the last
factor in the denominator is $|d(r_+\circ\gamma)/dt|$.  For the vertical
family $\gamma(t)=(r_0,t)$ this becomes
\begin{equation}
\label{eq:vertical-two-flux-special-case}
 \rho^{\rm raw}_{r_+}=
 \frac{|a_\lambda|J_{{\rm chart},\lambda}J_{{\rm hol},\lambda}}
 {|\nabla G_\lambda|}\frac{\cos\varphi_-\cos\varphi_+}{\tau},
\end{equation}
which is precisely the calculation of Lemma~5.1 of
\cite{CanestrariHoles}.  Formula \eqref{eq:radial-two-flux-density}, not the
vertical special case, is used for a general physical face.  The reversed
orientation has the same change-of-variables form.
\end{lemma}

\begin{proof}
Input arclength coarea and the collision--SRB density give the measure
\[
 b_\lambda(t)\,dt,
 \qquad b_\lambda(t)=
 \frac{|a_\lambda|J_{{\rm chart},\lambda}J_{{\rm hol},\lambda}
       \cos\varphi_-(t)}{|\nabla G_\lambda(\gamma(t))|}.
\]
The one-dimensional area formula for the injective oriented branch
$t\mapsto T\gamma(t)$ divides this density by the curve Jacobian
$J_\gamma T(t)=|DT(\gamma(t))\gamma'(t)|$, proving
\eqref{eq:radial-two-flux-density}.  Parametrizing the image by $r_+$ instead
divides by $|d(r_+\circ\gamma)/dt|$.  When $r_-$ is fixed, the billiard
differential identity
$dr_+/d\varphi_-=-\tau/\cos\varphi_+$ yields
\eqref{eq:vertical-two-flux-special-case}.  Applying the same area formula to
$T^{-1}$ proves the reversed statement without asserting that a general
tangent is vertical.
\end{proof}

\begin{definition}[Geometric flux and component condition
$\mathrm{FZ}_{\rm geom}$]
\label{def:flux-complexity}
A face family satisfies $\mathrm{FZ}_{\rm geom}(\theta_Z)$ if the density of
Lemma~\ref{lem:two-flux-coarea} has marks
$\mathfrak M_\lambda,\mathfrak C_\lambda\ge1$ such that
\begin{align}
 p_{\lambda,j,k,r}
 &\le C\mathfrak M_\lambda
 \frac{|W_{\lambda,j,k,r}|}{(1+j)^2(1+k)^2},
 \label{eq:radial-component-flux}\\
 \sum_{j,k,r}\frac1{(1+j)^2(1+k)^2}
 &\le\mathfrak C_\lambda:=CN_{\rm or}(\lambda),
 \label{eq:weighted-component-complexity}
\end{align}
and, for some $q_F>1$,
\begin{equation}
\label{eq:primitive-flux-ledger}
 \sum_\lambda M_\lambda
 \left[\operatorname{Mark}(\lambda)^2
 \left\{1+\frac{\mathfrak M_\lambda\mathfrak C_\lambda}
 {M_\lambda}\right\}\right]^{q_F}<\infty.
\end{equation}
Pieces whose tangent makes angle at most $\epsilon$ with both usable cones
have total marked mass at most $C\epsilon^{\theta_Z}$.  This is the required
near-stable angle-sublevel input.  On the complementary good-angle pieces,
the product of all inverse-angle costs in the graph chart, density,
nonstationary phase, and recovery maps is at most
$C\epsilon^{-A_{\rm ang}}$ for a fixed $A_{\rm ang}>0$.  Both constants are
part of $\mathrm{FZ}_{\rm geom}$.  This condition contains no multi-time
mixing or split-law assertion.
\end{definition}

\begin{definition}[Split-law standard-family and positive-envelope interface]
\label{def:split-law-envelope}
Let $F_{s,L}^{\rm col}$ be the sharp four-state functional of the assembled
finite-parameter collar, after coincident artificial faces have cancelled.
For every $(m,n)$, expansion of the global test-side double projection gives
the exact signed law
\begin{equation}
\label{eq:exact-split-law-expansion}
 \Sigma_{m,n}^{\rm law}
 =\rho_{m,n}^{0}-\sum_{j=1}^{J}a_j\rho_{m,n}^{j},
 \qquad \rho_{m,n}^{0}=\rho_{m,n}^{\rm orig},
\end{equation}
where the $\rho_{m,n}^{j}$ are the block-product laws in the functional-
correlation split; all $\rho_{m,n}^{j}$, including $\rho_{m,n}^{0}$, are
probability laws, $J\le J_0$, and
$Z_{\rm law}:=1+\sum_j|a_j|\le Z_0$ uniformly.  Define the positive envelope
\begin{equation}
\label{eq:positive-envelope-law}
 \bar\rho_{m,n}:=Z_{\rm law}^{-1}
 \left(\rho_{m,n}^{0}+\sum_{j=1}^{J}|a_j|\rho_{m,n}^{j}\right).
\end{equation}
The interface $\mathrm{SF}_{\rm split}$ requires that the restriction of every
constituent law in \eqref{eq:exact-split-law-expansion}, in either usable time
orientation, disintegrates into bounded-mark proper standard families with the
same conditional mark--$Z$ ledger as the original collar.  This is a primitive
joint-law statement, not a consequence of the one-coordinate marginals.

Let $\mathcal P_r^{m,n}$ be the common rank-$r$ product-cylinder partition,
refined into connected event-chart components.  For a bounded $F$ put
\begin{equation}
\label{eq:envelope-cylinder-oscillation}
 \Omega_r^{\rm env}(F):=
 \sum_{C\in\mathcal P_r^{m,n}}\bar\rho_{m,n}(C)
 \operatorname*{ess\,osc}_{C,\bar\rho_{m,n}}F.
\end{equation}
The connected-cylinder pullback condition $\mathrm{CP}_{\rm env}$ says that,
outside pieces of total envelope mass $Ce^{-c_{\rm bad}r}$, every such
component pulls back to a good-angle homogeneous event cell of physical
diameter at most $Ce^{-\chi r}$, with the conditional density and component
marks in $\mathrm{FZ}_{\rm geom}$ and $\mathrm{SF}_{\rm split}$.  Grazing,
near-stable, and homogeneity-boundary cells are included in the displayed bad
mass, with every component counted.
There are fixed $A_{\rm amp}>0$ and $\alpha_{\rm amp}\in(0,1]$ such that, on
a good cell, the smooth coefficient multiplying each sharp collar indicator
has oscillation at most
$C|s|^{-1}e^{A_{\rm amp}L}e^{-\chi\alpha_{\rm amp}r}$.
The connected refinement is required to be separation-time admissible: two
distinct retained atoms either receive the same cylinder value or differ in
some coordinate before rank $r$.  Hence a function constant on its atoms has
separate dynamic-H\"older cost at most $Ce^{q_{\rm cyl}r}$ times its
supremum.  In particular no artificial chart boundary may create an
uncharged jump between symbolically indistinguishable points.
\end{definition}

\begin{proposition}[Geometric envelope dynamical tube]
\label{prop:envelope-dynamical-tube}
Assume $\mathrm{FZ}_{\rm geom}(\theta_Z)$,
$\mathrm{SF}_{\rm split}$, and $\mathrm{CP}_{\rm env}$.  Then there are
$A_{\rm col}>0$ and $\alpha_{\rm col}\in(0,1]$ such that
\begin{equation}
\label{eq:FZ-dynamical-tube}
 \sup_{m,n}\Omega_r^{\rm env}(F_{s,L}^{\rm col})
 \le C|s|^{-1}e^{A_{\rm col}L}e^{-\alpha_{\rm col}r},
 \qquad r\ge1.
\end{equation}
\end{proposition}

\begin{proof}
On a good connected component, split the oscillation into variation of the
smooth coefficient and variation of a sharp indicator.  The first part is
bounded by the scaled H\"older clause in $\mathrm{CP}_{\rm env}$.  The second
produces two points on opposite sides of a genuine face.  Connectedness and
the diameter bound imply that the whole component lies in the
$Ce^{-\chi r}$ tube of that face.  By
\eqref{eq:positive-envelope-law}, its envelope mass is the fixed positive
combination of the actual and split-law masses.  Apply the full tangent area
formula and the standard-family conditional density bound to every
constituent.  The good-angle part costs
$C|s|^{-1}e^{A_{\rm col}L}e^{-\chi\alpha_0r}$ for a geometric tube exponent
$\alpha_0>0$, after taking
$A_{\rm col}\ge A_{\rm amp}$ and decreasing $\alpha_0$ to at most
$\alpha_{\rm amp}$.  The complement costs
$C|s|^{-1}e^{A_{\rm col}L}e^{-c_{\rm bad}r}$ by
$\mathrm{CP}_{\rm env}$.  Balancing the angle cutoff with the
$\epsilon^{\theta_Z}$ sublevel and decreasing the exponent once gives
\eqref{eq:FZ-dynamical-tube}.  Artificial boundaries were cancelled before
\eqref{eq:exact-split-law-expansion}, and connected components were indexed
before the diameter argument, so neither can create an uncharged crossing.
The standard-pair strip coverage of \cite{CanestrariPerturb} is a template
for verifying $\mathrm{CP}_{\rm env}$; it is not used as a black box.
\end{proof}

\begin{proposition}[Flux-weighted radial boundary lemma]
\label{prop:radial-direct-Z}
Assume $\mathrm{FZ}_{\rm geom}(\theta_Z)$ and decompose the signed amplitude into
positive pieces.  Then the normalized standard family obeys
\begin{equation}
\label{eq:radial-direct-Z-bound}
 M_\lambda Z(\mathcal G_\lambda)
 =\sum_{j,k,r}\frac{p_{\lambda,j,k,r}}
 {|W_{\lambda,j,k,r}|}
 \le C\mathfrak M_\lambda\mathfrak C_\lambda,
\end{equation}
and, for the actual physical labels,
\begin{equation}
\label{eq:joint-mark-Z-ledger}
 \sum_\lambda M_\lambda
 \operatorname{Mark}(\lambda)^2
 \{1+Z(\mathcal G_\lambda^-)+Z(\mathcal G_\lambda^+)\}<\infty.
\end{equation}
and the near-stable pieces have the marked
$\epsilon^{\theta_Z}$ tail in Definition~\ref{def:flux-complexity}.
\end{proposition}

\begin{proof}
Divide \eqref{eq:radial-component-flux} by the component length and sum over
the full triple $(j,k,r)$.  Equation
\eqref{eq:weighted-component-complexity} gives
\eqref{eq:radial-direct-Z-bound}; no component is suppressed.  Multiply by
the squared density/amplitude mark and use
\eqref{eq:primitive-flux-ledger} (first with power one) to obtain
\eqref{eq:joint-mark-Z-ledger}.  The full tangent Jacobian, weighted
component complexity, and angle sublevel are the three falsifiable geometric
inputs; none is deduced merely from the homogeneity labels.
\end{proof}

\begin{lemma}[Explicit angle balance]
\label{lem:FZ-angle-balance}
Suppose the good-angle core of a central or off-diagonal estimate is
$Ce^{-c_{\rm core}N}$ before inverse-angle loss.  Set
\[
 \epsilon_N=e^{-\zeta N},\qquad
 \zeta=\frac{c_{\rm core}}{A_{\rm ang}+\theta_Z}.
\]
Then the sum of the bad-angle and good-angle contributions is at most
$Ce^{-c_ZN}$, where
\begin{equation}
\label{eq:explicit-cZ}
 c_Z=\frac{c_{\rm core}\theta_Z}
 {A_{\rm ang}+\theta_Z}>0.
\end{equation}
\end{lemma}

\begin{proof}
The bad-angle mass is $e^{-\zeta\theta_ZN}$, while the good-angle estimate is
$e^{-(c_{\rm core}-A_{\rm ang}\zeta)N}$.  The displayed choice makes the two
exponents equal.
\end{proof}

\begin{corollary}[Uniform radial recovery]
\label{cor:radial-bi-recovery}
For each time orientation, the radial family is regular after at most
\begin{equation}
\label{eq:radial-log-Z-recovery}
 R^\pm\le \alpha_1\log\{1+Z(\mathcal G_{\rm rad}^\pm)\}+\alpha_2
\end{equation}
iterates.  After multiplication by physical mass, the recovery-window sum in
all four time quadrants is finite and is controlled by
\eqref{eq:joint-mark-Z-ledger}.
\end{corollary}

\begin{proof}
Apply \cite[Corollary~6.15(c)]{CanestrariHoles} to the positive standard
families in each orientation.  Use
$\log(1+Z)\le C_\varepsilon(1+Z)^\varepsilon$ and multiply by
$M_\lambda$.  The joint mark--$Z$ ledger then sums both orientations and all
pre-recovery iterates.  This is the weighted hypothesis of
Proposition~\ref{prop:recovery-window-summability}.
\end{proof}

\begin{lemma}[Observable multiplier to positive proper families]
\label{lem:signed-observable-multiplier}
Let $(W,\rho)$ be a positive standard pair and let
$q_m=(g\circ T^{-m})\rho$ on the time orientation for which $W$ is unstable.
There is $C_m\le C\|g\|_{C^\alpha}$ such that
\[
 q_m=(q_m+C_m\rho)-C_m\rho
\]
is a difference of positive densities.  Their total masses are bounded by
$C\|g\|_{C^\alpha}$ and their scaled log--H\"older marks are at most
$Ce^{\chi m}\operatorname{Mark}(W,\rho)$.  For every admissible slowly
moving sequence in the compact table class, after
\begin{equation}
\label{eq:multiplier-recovery-time}
 R_m\le C_0+C_1m+C_2\log\{1+Z(\mathcal G)\}
\end{equation}
iterates, both positive families are proper.  The reversed collision map has
the analogous statement.
\end{lemma}

\begin{proof}
Take $C_m=2\|g\|_\infty$.  Positivity and the mass bound are immediate.
Hyperbolicity and bounded distortion give
$[g\circ T^{-m}]_{\alpha,\rm sc}\le Ce^{\chi m}\|g\|_{C^\alpha}$.
Since the added density is at least $C_m\rho/2$, division by it gives the same
scaled log--H\"older bound.  More explicitly, the density recursion in
Lemma~13 of \cite{SYZ} has the form
$H_{j+1}\le\theta H_j+C$ on every homogeneous descendant.  Starting from
$H_0\le Ce^{\chi m}\operatorname{Mark}(W,\rho)$, it reaches the fixed regular
cone after $C_0+C_1m$ iterates.  The normalized boundary recursion in
Lemma~16 of \cite{SYZ} is
$Z_j/M\le C\{1+\theta_p^j Z_0/M\}$, so a further
$C_2\log(1+Z_0/M)$ iterates make the family proper.  Theorems~1'--2' of
\cite{SYZ} make all constants uniform for the compact slowly moving table
class.  The fixed-map statement follows by taking a constant sequence, and
time reversal gives the last assertion.
\end{proof}

\begin{proposition}[Signed proper-family mapping and off-diagonal bounds]
\label{prop:signed-off-diagonal-mapping}
Assume the forward and reversed billiard growth/coupling theorem on positive
proper standard families, $\mathrm{FZ}_{\rm geom}(\theta_Z)$, and the scaled density
bound of Proposition~\ref{prop:scaled-face-density}.  Decompose every signed
face amplitude as the difference of two positive amplitudes before
normalization.  Then there are positive exponents
$\alpha_\pm,\beta_\pm$ such that the complete face current satisfies
\begin{align}
 |B_{m,n}|&\le Ce^{\beta_-m-\alpha_+n},
 \label{eq:verified-forward-cone}\\
 |B_{m,n}|&\le Ce^{-\alpha_-m+\beta_+n}.
 \label{eq:verified-backward-cone}
\end{align}
The constants include all recovery windows and are uniform under the
physical and flux ledgers.
\end{proposition}

\begin{proof}
Apply Lemma~\ref{lem:signed-observable-multiplier} on each positive face
piece.  Its recovery time is bounded by
\eqref{eq:multiplier-recovery-time}.  Equations
\eqref{eq:radial-direct-Z-bound} and
\eqref{eq:radial-log-Z-recovery} control the boundary part of that recovery.
Coupling for the remaining
$n-R^+_\lambda$ iterates gives
\[
 CM_\lambda\operatorname{Mark}(\lambda)^2(1+Z_\lambda^+)
 e^{\beta_-m-\alpha_+(n-R^+_\lambda)_+}.
\]
Before recovery use total variation.  The inequality
 $e^{\alpha_+R^+_\lambda}\le
 Ce^{\alpha_+C_1m}(1+Z_\lambda^+)^{\alpha_+C_2}$ and
\eqref{eq:joint-mark-Z-ledger}, with a weakened coupling exponent chosen below
the available mark--$Z$ moment, sum all labels and recovery windows.  Absorb
$e^{\alpha_+C_1m}$ into $e^{\beta_-m}$.  This is
\eqref{eq:verified-forward-cone}.  Apply the same argument to the reversed
collision map for \eqref{eq:verified-backward-cone}; positive/negative pieces
are recombined only after the two bounds are established.
\end{proof}

\section{Slope-energy and fractional-phase proof architecture}

The proof now has six noncircular links.
\begin{enumerate}[label=\textup{(F\arabic*)}]
\item Quotient the exponential Young return on the regular magnet to the
      one-state countable $C^{1+\alpha}$ IFS of Proposition
      \ref{prop:marked-full-branch-quotient}.  Physical side labels remain
      in the signed amplitude fiber.
\item Lemma~\ref{lem:homogeneity-grazing-moment} combines the exponential
      return tail with the maximal grazing rank and proves the joint moment
      needed for the countable quotient and the marked coarea amplitudes.
\item The normal-cycle value $\mathfrak A=-1$ motivates the clean template;
      $\mathrm{SS}$ and the finite-depth physical thickening $\mathrm{FT}$, together
      with the mass-preserving buffer, give the energy theorem and Jensen.
\item Proposition~\ref{prop:exact-reduced-slope} and smooth value-space
      cutoffs and Lemma~\ref{lem:fractional-nonstationary-phase} give the
      balanced Fourier estimate without curvature or componentwise phase
      partitions.
\item Global centering is performed before the physical chart sum.
      Lemma~\ref{lem:dominant-coordinate-reduction} removes contracted modes;
      Proposition~\ref{prop:functional-correlation-low-frequency} treats the
      remaining low modes, while Proposition
      \ref{prop:central-phase-verification} treats the high modes.
\item At each cutoff coarea currents converge at fixed depth.  The
      component-indexed $\mathrm{FZ}_{\rm geom}$ estimate, the
      $\mathrm{SF}_{\rm split}$ and $\mathrm{CP}_{\rm env}$ interfaces,
      $\mathrm{QT}_{\rm path}$, and the strict scalar-tail margin remove
      $L,H$ before recovery summation.
\end{enumerate}

This order makes clear which external results are used.  Leclerc's
Lemma~3.15 supplies the cohomological bracket and his
Propositions~4.31--4.33 supply the four-point all-scale polynomial escape.
The final oscillatory estimate is the elementary fractional phase lemma
proved above; LPS is retained only as background for the quotient route.
Bertozzi's triple inducing supplies a useful model for the one-state gate
architecture, not the slope conclusion.  The positive entrance/exit comparison is the symbolic local
product argument of \cite{ClimenhagaLPS}.  The
one-step expansion, linear singularity complexity, and uniform billiard
constants used by the master atlas are summarized in \cite{DLreview}.

The route deliberately does not use the recent curvilinear two-ends theorem
or $L^2$ flattening: their nonlinear applications require $C^2$ or higher
geometry, whereas the billiard quotient available here is only
$C^{1+\alpha}$.  It also does not infer a product tower tail from two
marginal Chernoff bounds: the diagonal product exponent comes from time
comparability and functional correlation, while the off-diagonal regions are
handled before the three-region assembly.

\subsection{Execution audit}

The central analytic chain is now represented by explicit propositions.
The remaining primitive geometric inputs are explicitly named
$\mathrm{SS}$, $\mathrm{FT}$, $\mathrm{BR}_{\rm sec}$,
$\mathrm{FZ}_{\rm geom}$,
$\mathrm{SF}_{\rm split}$, $\mathrm{CP}_{\rm env}$, and pathwise
$\mathrm{QT}$.  The exact gate interface is
Lemma~\ref{lem:cohomological-bracket-identification}, Proposition
\ref{prop:context-polynomial-blowup}, and Lemma~\ref{lem:gate-buffer}; the
slope-energy theorem is Theorem~\ref{thm:direct-slope-small-ball}; the
balanced phase is Proposition~\ref{prop:central-phase-verification}; and the
low-frequency collision-time estimate is Proposition
\ref{prop:functional-correlation-low-frequency}.  Radial recovery is applied
only after these estimates through Corollary~\ref{cor:radial-bi-recovery} and
Proposition~\ref{prop:recovery-window-summability}.

The finite-parameter collar is treated below by the actual/split-law cylinder
approximation theorem; no fixed-slice smoothing estimate is used.

\begin{lemma}[Parameter derivatives of contracting return words]
\label{lem:parametric-inverse-word}
Let $h_{j,s}$ be the inverse one-step branches in a fixed compact regular
atlas and assume, uniformly in $j$ and $|s|<s_0$,
\[
 \|Dh_{j,s}\|_\infty\le\lambda<1,
 \qquad
 \|\partial_sh_{j,s}\|_{C^1}
 +\|\partial_s\log|h_{j,s}'|\|_\infty\le A.
\]
For a clean word $a=(a_1,\ldots,a_R)$ let
$h_{a,s}=h_{a_1,s}\circ\cdots\circ h_{a_R,s}$.  Then
\begin{equation}
\label{eq:parametric-word-bounds}
 \|\partial_s h_{a,s}\|_\infty\le \frac{A}{1-\lambda},
 \qquad
 \|\partial_s\log|h_{a,s}'|\|_\infty\le C(1+R).
\end{equation}
The same linear-in-$R$ bound holds for products of uniformly $C^1$
Jacobian, stable-holonomy, and coarea factors.  Hence an exponential return
moment dominates the difference quotients of all clean marked amplitudes and
their return tail.
\end{lemma}

\begin{proof}
Differentiate the composition.  The contribution created at position $k$ is
multiplied on the left by at least $k-1$ contracting derivatives, so its
$C^0$ norm is bounded by $A\lambda^{k-1}$; summation gives the first bound.
The logarithmic derivative of a composition is a sum of $R$ one-step
logarithmic derivatives.  Differentiating it gives a uniformly bounded direct
term at every position and orbit-displacement terms controlled by the first
bound and uniform $C^2$ distortion.  This gives $C(1+R)$.  The product
statement follows after taking logarithms.  Finally
$R^p e^{-\epsilon R}$ is summable for every fixed $p$, so the exponential
return moment absorbs all these factors.
\end{proof}

This lemma shows that long return words do not create an exponential
parameter-derivative loss on the inverse-branch side.  The remaining issue is
not the clean word calculus but the common refinement of words that approach
or cross a moving singularity; those pieces must be charged to the already
constructed face/recovery tree rather than differentiated as clean branches.

That singular refinement is most naturally made before fixing the parameter.

\subsection{Parameter-state geometry and cutoff removal}
\label{sec:parameter-state}

\begin{definition}[Pathwise quantitative jet-tail condition]
\label{def:quantitative-circular-jets}
Let $a(s)$ be the chosen analytic or circular table deformation with
$a(0)=a_0$.  For every face pair of depth at most $L$, let
$\Delta_{\lambda\lambda'}(s)$ be a finite collection of jet minors covering
the compact common-zero set after coincident germs are grouped.  Add the
individual face--cone boundary determinants and their critical-point jets;
these are the functions used to form the oriented tangent cuts in Lemma
\ref{lem:cone-monotone-components}.  The path
satisfies $\mathrm{QT}_{\rm path}$ if there are
$A,B,c_{\rm lab},A_{\rm loss},c_{\rm tail}>0$ such that
\begin{align}
 |\Delta_{\lambda\lambda'}(0)|&\ge e^{-AL},
 &\sup_{|s|<\bar s}|\partial_s\Delta_{\lambda\lambda'}(s)|&\le Ce^{BL},
 \label{eq:pathwise-minor-bounds}\\
 \#\{(\lambda,\lambda'):|\lambda|\le L\}&\le Ce^{c_{\rm lab}L}.
 \label{eq:pathwise-label-count}
\end{align}
All inverse-angle, coarea, chart, density, and parameter-derivative losses on
the depth-$L$ good set are at most $Ce^{A_{\rm loss}L}$, while the unsmoothed
return/grazing mass is at most $Ce^{-c_{\rm tail}L}$.  The strict path margin
is
\begin{equation}
\label{eq:strict-jet-tail-margin}
 c_{\rm tail}-A_{\rm loss}>A+B.
\end{equation}
For a finite-dimensional ambient family, the sublevel estimate
\[
 \operatorname{Leb}\{a:|\Delta_{\lambda\lambda'}(a)|<e^{-AL}\}
 \le Ce^{-A\tau L},\qquad A\tau>c_{\rm lab},
\]
is a sufficient Borel--Cantelli mechanism for the first inequality at almost
every base parameter, but it is not a substitute for the derivative bound or
the path margin.
\end{definition}

\begin{theorem}[Product-a.e. compact analytic realization of the pathwise jet bound]
\label{thm:analytic-QT-realization}
Let each scatterer boundary vary in a fixed complex analytic strip of width
$\sigma_a>0$.  Fix a strip slack $\delta_a>0$, put
$\sigma_p=\sigma_a+\delta_a$, and parametrize a neighborhood by
\begin{equation}
\label{eq:compact-analytic-product-probe}
 a_q=e^{-\sigma_p|q|}b_q,\qquad b_q\in[-1,1].
\end{equation}
Equip the $b_q$ with an independent compactly supported product law whose
one-dimensional densities are uniformly bounded.  The support of
\eqref{eq:compact-analytic-product-probe} is compact in the width-$\sigma_a$
analytic topology and hence in the chosen compact $C^3$ table chart.  Suppose
the depth-$L$ face atlas has at most $Ce^{c_{\rm lab}L}$ noncoincident jet
charts and that
\begin{equation}
\label{eq:localized-analytic-jet-direction}
 \begin{gathered}
 \text{each compact jet chart has a cover by at most }
 Ce^{c_{\rm zero}L}\text{ fiber boxes;}\\
 \text{on every box one index }q\le C_0L\text{ is fixed for the }b_q
 \text{ restriction;}\\
 \text{that restriction has at most }Ce^{c_{\rm zero}L}
 \text{ monotonicity intervals, and}\\
 |\partial_{b_q}\Delta|\ge
 e^{-(B_0+\sigma_pC_0)L}
 \quad\text{whenever }|\Delta|\le2e^{-AL}.
 \end{gathered}
\end{equation}
The fiber condition may equivalently be supplied by a uniform
Remez--Cartan/valency estimate with the same exponent $c_{\rm zero}$.  Then,
for every
\begin{equation}
\label{eq:analytic-QT-window}
 B_0+\sigma_pC_0+c_{\rm lab}+2c_{\rm zero}
 <A<c_{\rm tail}-A_{\rm loss}-B,
\end{equation}
product-almost every analytic base table satisfies
$|\Delta_{\lambda\lambda'}(0)|\ge e^{-AL}$ for all sufficiently large $L$.
For every analytic direction whose depth-$L$ derivative cost is at most
$e^{BL}$, the resulting path satisfies $\mathrm{QT}_{\rm path}$.  Small
analytic perturbations remain in the finite-horizon dispersing class.

For a fixed finite branch family, the adapted boundary-perturbation
construction of \cite{BFLS2026} supplies the relevant finite-jet
surjectivity template: perturb a boundary arc at the earliest differing
collision and then realize the required finite jet by an analytic
trigonometric polynomial.  Condition
\eqref{eq:localized-analytic-jet-direction} is the quantitative, simultaneous
version needed here.  Finite-branch jet surjectivity alone does not control
fiber valency.  For analytic minors one may verify the stated condition by a
quantitative propagation-of-smallness/Remez estimate such as
\cite{HWWRemez}; if the branch equations form a Nash family, the
Bernstein--Remez theorem \cite{BarbieriNash} is an alternative.  The theorem
applies only when the corresponding strip, valency, and normalized-coordinate
costs are verified for every face and cone-angle jet in the master atlas.  It
makes no assertion for the finite-dimensional circular subfamily.
\end{theorem}

\begin{proof}
Condition on all normalized Fourier coefficients except $b_q$ furnished
by \eqref{eq:localized-analytic-jet-direction}.  The one-dimensional change
of variables on each monotonicity interval and the bounded conditional density
give
\[
 \mathbb P\{|\Delta|<e^{-AL}\}
 \le C e^{2c_{\rm zero}L}
 e^{-\{A-B_0-\sigma_pC_0\}L}.
\]
The first $e^{c_{\rm zero}L}$ counts fiber boxes and the second counts
monotonicity intervals; the exponential Fourier weight is already present in
$\partial_{b_q}$.  Union over at most $Ce^{c_{\rm lab}L}$ labels.  The first inequality
in \eqref{eq:analytic-QT-window} makes the resulting series in $L$
summable, so Borel--Cantelli gives the base-minor bound.  Analytic strip norms
control the derivative along a fixed analytic direction by $Ce^{BL}$.
The second inequality in \eqref{eq:analytic-QT-window} is exactly the strict
path margin.  Finally, strict convexity, disjointness, and finite horizon are
open in the chosen compact $C^3$ neighborhood.
\end{proof}

\begin{lemma}[Persistence up to a depth depending on the parameter]
\label{lem:finite-depth-face-separation}
Choose
\begin{equation}
\label{eq:kappa-path-window}
 \frac1{c_{\rm tail}-A_{\rm loss}}<\kappa<\frac1{A+B},
 \qquad L(s)=\lfloor\kappa\log(1/|s|)\rfloor.
\end{equation}
For all sufficiently small $s$ and all $L\le L(s)$,
\[
 |\Delta_{\lambda\lambda'}(\theta s)|\ge\tfrac12e^{-AL},
 \qquad 0\le\theta\le1.
\]
Thus every noncoincident depth-$L(s)$ face family remains transverse along the
whole coarea segment from $0$ to $s$.
\end{lemma}

\begin{proof}
The mean-value theorem and \eqref{eq:pathwise-minor-bounds} give
\[
 |\Delta(\theta s)-\Delta(0)|\le C|s|e^{BL}.
\]
For $L\le L(s)$ the ratio of this error to $e^{-AL}$ is at most
$C|s|^{1-\kappa(A+B)}=o(1)$.  The interval for $\kappa$ is nonempty precisely
because of \eqref{eq:strict-jet-tail-margin}.
\end{proof}

\begin{proposition}[Depth-dependent coarea and pathwise scalar tail]
\label{prop:parameter-state-master-refinement}
Assume $|\nabla_xG_\lambda|\ge c_G>0$ on every active face.  For all depths
$L\le L(s)$,
\begin{equation}
\label{eq:master-coarea-difference}
 \frac{\one_{\{G_\lambda(s,\cdot)>0\}}
             -\one_{\{G_\lambda(0,\cdot)>0\}}}{s}
 =\int_0^1\delta(G_\lambda(\theta s,\cdot))
   \partial_sG_\lambda(\theta s,\cdot)\,d\theta
\end{equation}
is a uniformly controlled coarea current.  The omitted deeper words satisfy
\begin{equation}
\label{eq:finite-depth-tail}
 |s|^{-1}\,Ce^{-(c_{\rm tail}-A_{\rm loss})L(s)}
 \le C|s|^{\kappa(c_{\rm tail}-A_{\rm loss})-1}=o(1).
\end{equation}
This is a scalar total-variation tail.  Its conversion to an $o(1)$ tail in
the infinite CM2 double sum is a separate time--depth argument below; no such
conversion is inferred from \eqref{eq:finite-depth-tail} alone.  No fixed
neighborhood is asserted to be uniformly transverse at every depth.
\end{proposition}

\begin{proof}
Lemma~\ref{lem:finite-depth-face-separation} gives the coarea identity and all
geometric bounds along $\theta s$ for shallow words.  The return/grazing tail,
the joint mark--$Z$ ledger, and the allowed geometric loss give
$Ce^{-(c_{\rm tail}-A_{\rm loss})L}$ before the difference quotient.  Apply
the total-variation bound to both endpoints for words deeper than $L(s)$ and
divide by $|s|$.  The first inequality in
\eqref{eq:kappa-path-window} gives \eqref{eq:finite-depth-tail}.
\end{proof}

\begin{lemma}[Finite-parameter collar as a law-adapted functional]
\label{lem:finite-s-collar-admissibility}
Assume $\mathrm{FZ}_{\rm geom}(\theta_Z)$,
$\mathrm{SF}_{\rm split}$, and $\mathrm{CP}_{\rm env}$.
Let $\mathcal R_s^{>L}$ be the complete double-centered part of
$s^{-1}(P_s-P_0)$ carried by return/grazing words deeper than $L$.  Before
the coarea limit it has a full-dimensional density in every local event chart,
but its CM2 pairing is
\begin{equation}
\label{eq:finite-s-collar-law-pairing}
 B^{\rm deep}_{m,n}(s)=
 \int F_{s,L}^{\rm col}\,d\Sigma_{m,n}^{\rm law},
\end{equation}
where $\Sigma_{m,n}^{\rm law}$ is the exact finite expansion
\eqref{eq:exact-split-law-expansion}.  It is therefore not an integral against
$\mu^4$.
There are fixed $A_{\rm col}>0$ and $\alpha_{\rm col}\in(0,1]$ such that
\begin{align}
 \|\mathcal R_s^{>L}\|_{\rm TV}
 &\le C|s|^{-1}e^{-(c_{\rm tail}-A_{\rm loss})L},
 \label{eq:deep-collar-TV}\\
 \|F_{s,L}^{\rm col}\|_\infty
 &\le C|s|^{-1}e^{A_{\rm col}L},\qquad
 \Omega_r^{\rm env}(F_{s,L}^{\rm col})
 \le C|s|^{-1}e^{A_{\rm col}L}e^{-\alpha_{\rm col}r}.
 \label{eq:deep-collar-dynamical-tube}
\end{align}
The same pieces, disintegrated in either time orientation, are signed
differences of regular measured unstable families after the recovery time of
Lemma~\ref{lem:signed-observable-multiplier}; their total mass is bounded by
the right side of \eqref{eq:deep-collar-TV}.
\end{lemma}

\begin{proof}
In fixed transported coordinates the difference of the two moving domains is
a collar of width $O(|s|)$ and division by $s$ costs exactly one power of
$|s|^{-1}$.  Expanding the test-side projection before taking an absolute
value gives \eqref{eq:finite-s-collar-law-pairing} and
\eqref{eq:exact-split-law-expansion}.  The sharp indicator is not dynamically
H\"older.  Apply instead Proposition~\ref{prop:envelope-dynamical-tube}; this gives
\eqref{eq:deep-collar-dynamical-tube}.  Absolute summation with the
return/grazing mass gives \eqref{eq:deep-collar-TV}.  Disintegrating the collar
along the usable unstable orientation preserves total mass.  The density and
$Z$ recursions of Lemmas~13 and 16 and Theorems~1'--2' of \cite{SYZ}, together
with Lemma~\ref{lem:signed-observable-multiplier}, give the last assertion
uniformly for the moving table sequence.
\end{proof}

\begin{theorem}[Positive-envelope cylinder approximation]
\label{thm:law-adapted-cylinder-FCB}
Let $F$ be a bounded four-state functional and let the signed law and positive
envelope be \eqref{eq:exact-split-law-expansion}--
\eqref{eq:positive-envelope-law}.  Suppose
$\Omega_r^{\rm env}(F)\le Ve^{-\alpha r}$ and the common connected partition
is separation-time admissible as in $\mathrm{CP}_{\rm env}$.  On every rank-$r$ connected
product cylinder choose a value in the
$\bar\rho_{m,n}$-essential range of $F$ and call the resulting common cylinder
function $F_r$.  There is a fixed $q_{\rm cyl}>0$ such that
\begin{align}
 \|F_r\|_{\rm SH}&\le C\|F\|_\infty e^{q_{\rm cyl}r},
 \label{eq:cylinder-approximant-SH}\\
 \left|\int(F-F_r)\,d\Sigma_{m,n}^{\rm law}\right|
 &\le Z_0Ve^{-\alpha r}.
 \label{eq:cylinder-approximant-error}
\end{align}
Consequently a functional-correlation split across a gap $N$ obeys
\begin{equation}
\label{eq:cylinder-functional-correlation}
 |\operatorname{FC}_N(F)|
 \le C\left\{Ve^{-\alpha r}
       +\|F\|_\infty e^{-c_{\rm FC}N+q_{\rm cyl}r}\right\}.
\end{equation}
If two mutually singular constituent laws see the values zero and one on the
same cylinder, their positive envelope sees both values and the envelope
oscillation is one.  For the diagonal example, a value constant on every
constituent support is represented exactly; if a split constituent sees the
opposite side, the envelope charges that discrepancy.  Thus the construction
never replaces a law-adapted zero error by coordinatewise mollification across
the diagonal.
\end{theorem}

\begin{proof}
On each cylinder, the envelope-essential difference from the chosen value is
bounded by the envelope-essential oscillation.  Using
\eqref{eq:positive-envelope-law},
\[
 \int|F-F_r|\,d\rho^0+
 \sum_j|a_j|\int|F-F_r|\,d\rho^j
 =Z_{\rm law}\int|F-F_r|\,d\bar\rho_{m,n},
\]
which proves \eqref{eq:cylinder-approximant-error}.  Separation-time
admissibility in $\mathrm{CP}_{\rm env}$ says that a first possible jump
occurs only after some separation time drops below $r$; it costs at most
$C\theta_0^{-r}=Ce^{q_{\rm cyl}r}$ in that slot.  This proves
\eqref{eq:cylinder-approximant-SH}.  Apply \cite[Theorems~2.3--2.4]{LS} to
$F_r$ for the exact signed split and use
\eqref{eq:cylinder-approximant-error}.  The same $F_r$ serves all constituent
laws because it was constructed from their common positive envelope.  No
assertion about conditional coordinate marginals is used.
\end{proof}

\begin{theorem}[$(L,T,r)$ removal of the deep tail]
\label{thm:time-depth-diagonal}
Put $c_d=c_{\rm tail}-A_{\rm loss}$ and choose
$L(s)=\lfloor\kappa_L\log(1/|s|)\rfloor$ with
$\kappa_Lc_d>1$ and $\kappa_L(A+B)<1$, as in
\eqref{eq:kappa-path-window}.  Set
\begin{equation}
\label{eq:time-depth-threshold}
 c_{\rm col}:=
 \frac{\alpha_{\rm col}c_{\rm FC}}
      {q_{\rm cyl}+\alpha_{\rm col}},\qquad
 T(s)=\lfloor\kappa_T\log(1/|s|)\rfloor,\qquad
 \kappa_Tc_{\rm col}>1+\kappa_LA_{\rm col}.
\end{equation}
Let
\[
 B^{\rm deep}_{m,n}(s)=
 \langle Q_s^n\mathcal R_s^{>L(s)}Q_s^mg,f\rangle
\]
with the transported centered operators; replacing one side by $Q_0$ gives
the same conclusion.  Then
\begin{equation}
\label{eq:deep-tail-CM2-sum}
 \sum_{m,n\ge0}|B^{\rm deep}_{m,n}(s)|=o(1),
 \qquad s\to0.
\end{equation}
\end{theorem}

\begin{proof}
Write
$\delta_s=C|s|^{-1}e^{-c_dL(s)}
=O(|s|^{\kappa_Lc_d-1})=o(1)$.
For $N=\max\{m,n\}\le T(s)$, total variation and Markov normalization give
\[
 \sum_{m,n\le T(s)}|B^{\rm deep}_{m,n}(s)|
 \le C\delta_sT(s)^2=o(1).
\]
In the two off-diagonal regions, disintegrate the collar as in Lemma
\ref{lem:finite-s-collar-admissibility}.  Forward and reverse moving-table
coupling give the two cone bounds
$C\delta_se^{\beta_-m-\alpha_+n}$ and
$C\delta_se^{-\alpha_-m+\beta_+n}$.  Choosing the cone slopes as in Lemma
\ref{lem:three-region} makes their complete off-diagonal sums $O(\delta_s)$.

It remains to consider the central band with $N>T(s)$.  The collar itself,
not its singular coarea limit, is the sharp four-state function
$F_{s,L(s)}^{\rm col}$.  It is not dynamically H\"older.  Apply
Theorem~\ref{thm:law-adapted-cylinder-FCB} with the $N$-dependent rank
\[
 r_N=\left\lfloor
 \frac{c_{\rm FC}N}{q_{\rm cyl}+\alpha_{\rm col}}
 \right\rfloor.
\]
Global double centering kills both product terms, and
\eqref{eq:deep-collar-dynamical-tube}--
\eqref{eq:cylinder-functional-correlation} give
\[
 |B^{\rm deep}_{m,n}(s)|
 \le C|s|^{-1-A_{\rm col}\kappa_L}e^{-c_{\rm col}N}.
\]
There are $O(N)$ pairs with maximum $N$, so
\[
 \sum_{\substack{m\asymp n\\\max(m,n)>T(s)}}
 |B^{\rm deep}_{m,n}(s)|
 \le C|s|^{-1-A_{\rm col}\kappa_L}
       \sum_{N>T(s)}Ne^{-c_{\rm col}N}=o(1)
\]
by \eqref{eq:time-depth-threshold}.  The three regions exhaust the quadrant,
proving \eqref{eq:deep-tail-CM2-sum}.
\end{proof}

This construction is why a differentiable parameter family of Young towers
is unnecessary.  The dynamical return alphabet is fixed on clean master
cells, while every atom which changes its itinerary is represented by the
physical coarea current that the response formula already contains.

The clean cohomological bracket and the master refinement can now be put on one ledger.  The
point of the next proposition is not to create a second symbolic dynamics.
It records, on the amplitude side, which already existing full branches and
recovery descendants contain a physical face piece.

\begin{lemma}[Measurable refinement and conditional moments]
\label{lem:measurable-refinement-moments}
Let $\mathcal F_0\subset\cdots\subset\mathcal F_5$ be the finite or countable
label filtration obtained successively from: the physical chart/return word;
the two endpoint contexts; the finite native stopping antichain; the two
homogeneity cuts and component index; and the time orientation/path-depth
label.  At every stage let $R_i(\lambda_{i-1},d\lambda_i)$ be the normalized
mass kernel of the refinement.  Assume
\begin{equation}
\label{eq:refinement-kernel-mass}
 R_i(\lambda_{i-1},\Lambda_i)=1
 \quad\text{for every surviving parent,}
\end{equation}
and, with $\mathfrak F$ denoting the mark--$Z$ factor in
\eqref{eq:primitive-flux-ledger},
\begin{align}
 \mathbb E[(1+H)^{q_H}\mid\mathcal F_1]&\le C,&
 \mathbb E[\mathfrak F^{q_F}\mid\mathcal F_2]&\le C,
 \label{eq:conditional-refinement-moments}\\
 \mathbb E[\mathfrak d^{A_{\rm sec}q_D}\mid\mathcal F_1]&\le C,&
 \mathbb E[e^{q_R\epsilon(|c^-|+|c^+|)}\mid\mathcal F_0]&\le C.
 \notag
\end{align}
If $q_H^{-1}+q_F^{-1}+q_D^{-1}<1$ and $q_R>1$ is chosen above the remaining
conjugate exponent, then every product of the four displayed marks is
integrable under the terminal label law.  Moreover all coarser marginals are
unchanged by the refinement.
\end{lemma}

\begin{proof}
Equation \eqref{eq:refinement-kernel-mass} and Ionescu--Tulcea construct one
terminal probability law and show by iterated conditional expectation that
each earlier marginal is preserved.  Apply conditional H\"older first at the
component stage with exponents $q_F$ and its conjugate, then at the native and
context stages with $q_H,q_D,q_R$.  The strict reciprocal margin leaves a
finite conjugate exponent for the exponential context factor.  This proves
integrability without multiplying moments taken under different label laws.
Coincident artificial faces are cancelled before $R_1$ is defined, so the
construction also preserves the signed physical marginal.
\end{proof}

\begin{proposition}[Global physical ledger]
\label{prop:global-labelled-ledger}
Assume the exponential return/context moments and the conditional refinement
estimates \eqref{eq:conditional-refinement-moments},
$\mathrm{SS}(d_Z)$, $\mathrm{FT}(d_Z)$,
$\mathrm{BR}_{\rm sec}$, $\mathrm{FZ}_{\rm geom}(\theta_Z)$,
$\mathrm{SF}_{\rm split}$, $\mathrm{CP}_{\rm env}$, and the pathwise jet condition of
Definition~\ref{def:quantitative-circular-jets}.  Let $q_H,q_F,q_D>1$ be the
available context, flux, and dominant-scale moment exponents, and assume
\begin{equation}
\label{eq:holder-conjugate-margin}
 q_H^{-1}+q_F^{-1}+q_D^{-1}<1.
\end{equation}
After artificial-face cancellation, the complete first defect has a disjoint
signed refinement
\begin{equation}
\label{eq:global-label}
 \lambda=(\chi,a,k,c^-,g,c^+;Q;j,k',r,\sigma),
\end{equation}
where $\chi$ is a physical chart, $a$ a return word, $k$ the insertion time,
$c^-,g,c^+$ the single-source gate data, $Q$ a native gate ball,
$(j,k',r)$ the incoming/outgoing strips and connected-component index, and
$\sigma$ the chosen time orientation.  The refinement does not promote
amplitude cuts to new inverse branches.  For some $a_H>0$ below the context
moment threshold,
\begin{equation}
\label{eq:global-ledger-sum}
 \sum_\lambda M_\lambda
 e^{\epsilon(|c^-|+|c^+|)}
 (1+H_\lambda)^{a_H}\mathfrak d_\lambda^{A_{\rm sec}}
 \operatorname{Mark}(\lambda)^2
 \{1+Z(\mathcal G_\lambda^-)+Z(\mathcal G_\lambda^+)\}<\infty.
\end{equation}
All constants in the functional-correlation, native-energy,
fractional-phase, return, grazing, and direct-$Z$ estimates are uniform on
the quantitative table family.
\end{proposition}

\begin{proof}
Use the normalized refinement kernels of
Lemma~\ref{lem:measurable-refinement-moments}.  The context and $H$ moments
come from the full-branch Gibbs quotient and exponential gate hitting;
Proposition~\ref{prop:radial-direct-Z} gives the component-indexed
mark--boundary conditional moment before the boundary-oscillation theorem is
applied; and the dominant chart hypothesis gives the conditional
$\mathfrak d^{A_{\rm sec}q_D}$ moment.  Return/grazing tails and Proposition
\ref{prop:parameter-state-master-refinement} provide the last normalized
kernel; the strict margin \eqref{eq:strict-jet-tail-margin} absorbs its
finite-depth loss.  Apply \eqref{eq:conditional-refinement-moments} and the
strict conjugate margin \eqref{eq:holder-conjugate-margin}.  No unrecorded
$p>2$ moment and no intersection of unrelated label laws is used.  This proves
\eqref{eq:global-ledger-sum}; it is not assumed again in the main theorem.
\end{proof}

\begin{theorem}[Conditional pathwise quantitative CM2 criterion]
\label{thm:global-cm2-assembly}
Let an analytic-boundary or circular deformation $a(s)$ satisfy
$\mathrm{QT}_{\rm path}$.
Assume a compact finite-horizon chart with:
\begin{enumerate}[label=\textup{(H\arabic*)}]
\item the finite-horizon functional-correlation bound of \cite{LS}, the
      standard one-state marked quotient, and exponential return/context
      moments; a dominant inverse-scale moment
      $\sum_\lambda M_\lambda
      \mathfrak d_\lambda^{A_{\rm sec}q_D}<\infty$; and
      exponents $q_H,q_F,q_D$ satisfying
      \eqref{eq:holder-conjugate-margin}; all three moments hold conditionally
      along the common refinement as in
      \eqref{eq:conditional-refinement-moments};
\item forward/reverse proper-family growth and coupling uniformly for the
      compact slowly moving table class, in the form of Theorems~1'--2' of
      \cite{SYZ}; Lemma~\ref{lem:signed-observable-multiplier} then supplies
      the signed multiplier bridge;
\item $\mathrm{SS}(d_Z)$ and
      $\mathrm{FT}(d_Z)$ on the persistent clean gate, together with the
      dominant chart bounds \eqref{eq:primitive-dominant-chart-bounds} and
      the branch-section interface $\mathrm{BR}_{\rm sec}$;
\item the two-chart nondegeneracy $|\nabla G|\ge c_G$ and
      $\mathrm{FZ}_{\rm geom}(\theta_Z)$, including the full tangent coarea identity,
      weighted component complexity, angle sublevel, and good-angle exponent
      $A_{\rm ang}$, together with $\mathrm{SF}_{\rm split}$ and
      $\mathrm{CP}_{\rm env}$; Proposition~\ref{prop:envelope-dynamical-tube}
      then supplies \eqref{eq:FZ-dynamical-tube};
\item the regular defect satisfies
      $\sum_\chi\|\zeta_\chi k_{\rm reg}\|_{W^{R,1}}<\infty$ for some
      $1\le R\le3$, and the local face amplitudes have the scaled
      $\mathcal A^\alpha$ regularity used in
      Proposition~\ref{prop:scaled-face-density}; observable regularity is
      above the dominant two-dimensional Fourier threshold.  The joint face
      part of \eqref{eq:event-physical-norm} is concluded from
      Proposition~\ref{prop:global-labelled-ledger}, not assumed here.
\end{enumerate}
Then the complete double-centered physical first defect satisfies
\[
 |\langle Q^nK_VQ^mg,f\rangle|
 \le C\tau^m\vartheta^n\|g\|_{C^r}\|f\|_{C^r}.
\]
The assembled difference quotients converge after the depth-dependent split
$L(s)$ of \eqref{eq:kappa-path-window}; their deep remainders converge to zero
in the absolute CM2 double-sum norm by Theorem
\ref{thm:time-depth-diagonal}, so the CM2 series converges term by term as
$s\to0$.  More
precisely, if $c_{\rm off}$ is the off-diagonal coupling exponent and $c_Z$
is the exponent obtained by balancing the $\mathrm{FZ}_{\rm geom}$ angle sublevel, one
may choose
\begin{equation}
\label{eq:final-CM2-exponent}
 c_*<\min\{c_{\rm off},c_{\rm low},
 \eta\beta_{\rm cen},\eta R,
 c_{\rm tail}-A_{\rm loss},c_Z\}
\end{equation}
and then take $\tau=\vartheta=e^{-c_*/2}$ after changing the constant.
\end{theorem}

\begin{proof}
Proposition~\ref{prop:global-labelled-ledger} first supplies the complete face
part of \eqref{eq:event-physical-norm} from the primitive moment hypotheses.
Conditions $\mathrm{SS}$ and $\mathrm{FT}$, together with the
mass-preserving buffer, give
Theorem~\ref{thm:direct-slope-small-ball}.  The high-frequency balanced
estimate is Proposition~\ref{prop:central-phase-verification}.  The regular
two-dimensional defect is included in the low-frequency Proposition
\ref{prop:functional-correlation-low-frequency}.  Choosing
$K=e^{\eta\max\{m,n\}}$ produces the low, face-high, and regular-tail
exponents in
\eqref{eq:final-CM2-exponent}.

Outside the diagonal band use Proposition
\ref{prop:signed-off-diagonal-mapping};
Lemma~\ref{lem:three-region} converts the three estimates to product decay.
For radial pieces Proposition~\ref{prop:radial-direct-Z}, Lemma
\ref{lem:FZ-angle-balance}, and Corollary
\ref{cor:radial-bi-recovery} give the common physical majorant without a
component-length estimate and with every component explicitly summed.
Proposition~\ref{prop:global-labelled-ledger} gives the global physical
ledger from the primitive hypotheses.  Lemma
\ref{lem:finite-depth-face-separation} and Proposition
\ref{prop:parameter-state-master-refinement} give coarea convergence up to
$L(s)$.  Lemma~\ref{lem:finite-s-collar-admissibility} and Theorem
\ref{thm:time-depth-diagonal} give the required $o(1)$ CM2-norm tail beyond
$L(s)$.  No fixed neighborhood is
claimed to be transverse at every depth, and no qualitative residual
assertion is used in this step.
\end{proof}

\begin{corollary}[Product-a.e. analytic paths and the circular boundary]
\label{cor:analytic-pathwise-CM2}
Assume all hyperbolic, single-source, finite-thickening, branch-section,
geometric-flux, split-standard-family, and connected-pullback hypotheses of
Theorem~\ref{thm:global-cm2-assembly}.  If the analytic boundary family
satisfies the localized jet-direction condition
\eqref{eq:localized-analytic-jet-direction} and the exponent window
\eqref{eq:analytic-QT-window}, then for almost every base table in the
analytic product model and every admissible analytic direction with the
recorded derivative exponent, CM2 holds along the path.
For the finite-dimensional circular family the same conclusion holds under
the explicitly stated $\mathrm{QT}_{\rm path}$ condition; no automatic
circular transversality is asserted.
\end{corollary}

\begin{proof}
Apply Theorem~\ref{thm:analytic-QT-realization} and then Theorem
\ref{thm:global-cm2-assembly}.  The last sentence is the circular branch of
the latter theorem.
\end{proof}

\begin{proposition}[Pathwise coarea-current convergence]
\label{prop:common-refinement-majorant}
In a compact table chart satisfying \textup{(H1)--(H5)}, all constants in the
functional-correlation, native-energy, fractional-phase, direct-$Z$, and
parameter-tail estimates are uniform up to $L(s)$.  The shallow coarea current
has the CM2 majorant of Theorem~\ref{thm:global-cm2-assembly}, while the deep
remainder is $o(1)$ in the CM2 sum.
\end{proposition}

\begin{proof}
Hyperbolicity, distortion, the clean gate, and single-face nondegeneracy are
uniform on the quantitative compact atlas.  The
functional-correlation constants are uniform for a compact finite-horizon
family, and the return/context moments and both one-sided coupling estimates
are uniform by \textup{(H1)--(H2)}.  Proposition
\ref{prop:radial-direct-Z} fixes the component-indexed density/boundary
majorant.  Use the strict path margin
\eqref{eq:strict-jet-tail-margin}.  Proposition
\ref{prop:parameter-state-master-refinement} supplies the
depth-dependent scalar quotient tail.  Apply the low/high central estimates
and the off-diagonal couplings label by label to the shallow part, sum
\eqref{eq:global-ledger-sum}, and use Theorem
\ref{thm:time-depth-diagonal} for the deep part.
\end{proof}

The construction does not require a differentiable parameter family of
countable towers.  The coarea depth grows like $L(s)$, while the remaining
collision/grazing words are removed in the CM2 norm by the time--depth split
of Theorem~\ref{thm:time-depth-diagonal}, not by scalar total variation alone.

\section{Global ledger and quantitative falsification checks}

For reference, the logical dependencies of the active proof are
\[
\begin{array}{ccl}
 \textup{(H1)--(H3),(H5)}&\Longrightarrow&
   \text{branch-section finite-rank low + regular high},\\
 \textup{(H3),(H5)}&\Longrightarrow&
   \text{weighted energy + balanced/unbalanced phase},\\
 \textup{(H2),(H4)}&\Longrightarrow&
   \text{full-tangent }Z\text{ + two off-diagonal cones},\\
 \textup{(H1),(H2),(H4)},\mathrm{QT}_{\rm path}&\Longrightarrow&
   \text{envelope crossing + }(L,T,r)\text{ deep-tail removal},\\
 \textup{(H4)},\mathrm{QT}_{\rm path}&\Longrightarrow&
   \text{shallow coarea convergence}.
\end{array}
\]
The five outputs are combined only in
Theorem~\ref{thm:global-cm2-assembly}; none is used to define an input norm.

The proof has six independent exponent ledgers.
\begin{enumerate}[label=\textup{(A\arabic*)}]
\item In the diagonal band the low-frequency ledger is
\[
K^{A_{\rm fr}}e^{-c_-m-c_+n},
 \qquad K=e^{\eta\max\{m,n\}}.
\]
The branch-section contraction error
$e^{-c\max\{m,n\}}$ is charged to the same ledger before the cutoff.
\item The high-frequency ledger is
$K^{-\beta_{\rm cen}}$, obtained from the native single-source energy and
fractional phase.
\item The regular high-frequency ledger is $K^{-R}$ from
Lemma~\ref{lem:regular-high-tail}.
\item Lemma~\ref{lem:three-region} supplies the two off-diagonal coupling
ledgers.
\item The physical majorant is the pathwise return/grazing quotient tail
multiplied by the joint mark--$Z$ sum \eqref{eq:joint-mark-Z-ledger}.
\item The deep finite-parameter collar is removed in the absolute CM2 norm by
the time--depth ledger
\[
 |s|^{\kappa_L(c_{\rm tail}-A_{\rm loss})-1}\log^2(1/|s|)
 \quad\hbox{and}\quad
 |s|^{-1-A_{\rm col}\kappa_L}e^{-c_{\rm col}T(s)},
\]
whose two exponents are positive by
\eqref{eq:kappa-path-window} and \eqref{eq:time-depth-threshold}.
\end{enumerate}
Choosing $\eta$ as in Proposition
\ref{prop:functional-correlation-low-frequency} makes the first five
exponents positive, while the time--depth inequalities make the sixth ledger
vanish, and gives \eqref{eq:final-CM2-exponent}.  Zero and axis modes are
removed by global centering before the physical chart sum; they are not
estimated label by label.  The regular two-dimensional defect is included in
the same finite-rank low estimate and has its independent $W^{R,1}$ tail.

There are six falsifiable interfaces: $\mathrm{SS}$ must contain
exactly one Leclerc rough source; $\mathrm{FT}$ must thicken every finite
Leclerc child into an actual-SRB collision cylinder, with its weight
comparison and stopping-antichain cover treated as primitive physical inputs;
$\mathrm{BR}_{\rm sec}$ must produce same-branch homogeneous sections with
uniformly recoverable oscillatory densities; $\mathrm{FZ}_{\rm geom}$ must
count every connected component and prove the full tangent coarea Jacobian
(with the two-flux formula only in the vertical special case);
$\mathrm{SF}_{\rm split}$ must represent every actual and split-law block by
the stated bounded-mark standard families; $\mathrm{CP}_{\rm env}$ must give
connected event cells with a uniform diameter/tail estimate; and the chosen
deformation must verify fiberwise valency/Remez control and
$\mathrm{QT}_{\rm path}$ with the strict margin
\eqref{eq:strict-jet-tail-margin}.  These are primitive
quantitative conditions, not consequences hidden in a Banach norm or a
qualitative genericity statement.

\begingroup\raggedright
\begin{thebibliography}{99}
\bibitem{DLcones}
M.~Demers and C.~Liverani,
\emph{Projective cones for sequential dispersing billiards},
Commun. Math. Phys. 401 (2023), arXiv:2104.06947.

\bibitem{DLreview}
M.~Demers and C.~Liverani,
\emph{Recent progress in the application of transfer operators to dispersing
billiards}, arXiv:2606.10155 (2026).

\bibitem{DLsequential}
M.~Demers and C.~Liverani,
\emph{Central limit theorem for sequential dynamical systems},
arXiv:2502.07765v2 (2025).

\bibitem{SYZ}
M.~Stenlund, L.-S.~Young and H.-K.~Zhang,
\emph{Dispersing billiards with moving scatterers},
arXiv:1210.0011.

\bibitem{LS}
J.~Lepp\"anen and M.~Stenlund,
\emph{A note on the finite-dimensional distributions of dispersing billiard processes},
arXiv:1702.01461.

\bibitem{FlemingFCB}
N.~Fleming-V\'azquez,
\emph{Functional correlation bounds and optimal iterated moment bounds for
slowly-mixing nonuniformly hyperbolic maps}, arXiv:2106.06486v2.

\bibitem{BFLS2026}
I.~Baldom\'a, A.~Florio, M.~Leguil and T.~M.-Seara,
\emph{Chaoticity of generic analytic convex billiards},
arXiv:2605.19897 (2026).

\bibitem{Wormell}
C.~Wormell,
\emph{Conditional mixing in deterministic chaos},
arXiv:2206.09291; and
\emph{On convergence of linear response formulae in piecewise hyperbolic maps},
arXiv:2206.09292.

\bibitem{BB}
S.~Baker and A.~Banaji,
\emph{Polynomial Fourier decay for fractal measures and their pushforwards},
Math. Ann. 392 (2025), 209--261, arXiv:2401.01241.

\bibitem{BKS}
S.~Baker, O.~Khalil and T.~Sahlsten,
\emph{Fourier decay from $L^2$-flattening},
arXiv:2407.16699.

\bibitem{CD}
V.~Climenhaga and J.~Day,
\emph{Every finite horizon Sinai billiard map has a unique measure of maximal entropy},
arXiv:2604.25881.

\bibitem{BNST}
P.~B\'alint, P.~N\'andori, D.~Sz\'asz and I.~P.~T\'oth,
\emph{Equidistribution for standard pairs in planar dispersing billiard flows},
Ann. Henri Poincar\'e 19 (2018), arXiv:1707.03068.

\bibitem{HWWRemez}
S.~Huang, G.~Wang and M.~Wang,
\emph{Observability inequality, log-type Hausdorff content and heat equations},
arXiv:2411.11573v2 (2024).

\bibitem{BarbieriNash}
S.~Barbieri,
\emph{Bernstein--Remez inequality for Nash functions: a complex analytic
approach}, arXiv:2207.13922 (2022).

\bibitem{BCC}
O.~Butterley, G.~Canestrari and R.~Castorrini,
\emph{Discontinuities cause essential spectrum on surfaces},
Ann. Henri Poincar\'e (2024), arXiv:2306.00484.

\bibitem{CanestrariPerturb}
G.~Canestrari,
\emph{On linear response for discontinuous perturbations of smooth
endomorphisms}, Advances in Mathematics 497 (2026), 111008,
arXiv:2411.16628v3.

\bibitem{CanestrariHoles}
G.~Canestrari,
\emph{Linear response for Sinai billiards with small holes},
arXiv:2604.19671v2 (2026).

\bibitem{AlvesBahsoun}
J.~F.~Alves and W.~Bahsoun,
\emph{Linear response for skew-product maps with contracting fibres},
arXiv:2602.02317v2 (2026).

\bibitem{Leclerc}
G.~Leclerc,
\emph{Fourier decay of equilibrium states for bunched attractors},
arXiv:2301.10623.

\bibitem{Leclerc2026}
G.~Leclerc,
\emph{Fourier decay of equilibrium states and the Fibonacci Hamiltonian},
arXiv:2507.23731v2 (2026).

\bibitem{LPS}
G.~Leclerc, S.~Paukkonen and T.~Sahlsten,
\emph{Fourier decay in parabolic $C^{1+\alpha}$ systems with overlaps},
arXiv:2505.15468v3 (2025).

\bibitem{ARHW}
A.~Algom, F.~Rodriguez Hertz and Z.~Wang,
\emph{Polynomial Fourier decay and a cocycle version of Dolgopyat's method
for self conformal measures}, arXiv:2306.01275v2 (2024).

\bibitem{LOP}
Y.~Lima, D.~Obata and M.~Poletti,
\emph{Measures of maximal entropy for non-uniformly hyperbolic maps},
arXiv:2405.04676v2 (2025).

\bibitem{Chernov1999}
N.~Chernov,
\emph{Decay of correlations and dispersing billiards},
J. Stat. Phys. 94 (1999), 513--556.

\bibitem{Teixeira}
P.~Teixeira,
\emph{On the conservative pasting lemma},
Ergodic Theory Dynam. Systems 40 (2020), 1402--1440,
arXiv:1611.01694.

\bibitem{Ruehr}
R.~R\"uhr,
\emph{Pressure inequalities for Gibbs measures of countable Markov shifts},
arXiv:2012.13226v3.

\bibitem{Bertozzi}
N.~Bertozzi,
\emph{Exponential mixing for hyperbolic flows on non-compact spaces},
arXiv:2603.11852v2 (2026).

\bibitem{ClimenhagaLPS}
V.~Climenhaga,
\emph{Gibbs measures have local product structure},
arXiv:2310.17495v2 (2025).

\bibitem{GouezelRenewal}
S.~Gou\"ezel,
\emph{Berry--Esseen theorem and local limit theorem for non uniformly
expanding maps}, Ann. Inst. H. Poincar\'e Probab. Statist. 41 (2005),
997--1024, arXiv:math/0310391.

\bibitem{Galli}
D.~Galli,
\emph{A cohomological approach to Ruelle--Pollicott resonances and speed of
mixing of Anosov diffeomorphisms}, arXiv:2405.17045.
\end{thebibliography}
\endgroup

\end{document}
